% ===========================================================================
% supplement-main.tex — Supplementary Material master (the supplementary
% split of the main text).
%
% Contents (S-numbered sections):
%   S1  data provenance map            (sections/appendix-a-provenance.tex)
%   S2  tier-1 reproduction mechanics  (sections/appendix-b-reproducibility.tex)
%   S3  governing equations + frozen constants (sections/appendix-c-equations.tex)
%   S4  extended results, groups B and D (sections/supplement-d-extended-results.tex)
%   S5  full predicted-shift verification table
%                                      (sections/supplement-e-predicted-shifts.tex)
%   S6  video-artifact presentation methods (time remap / shutter)
%                                      (sections/supplement-f-timeremap.tex)
%
% CONVENTION: SIMPAT (Elsevier) accepts supplementary material as separate
% e-component files uploaded alongside the manuscript; this document is that
% file, built from the same section sources and the same references.bib as
% main.tex.
%
% CROSS-REFERENCES: \ref{...} into the MAIN text (sections, tables, figures
% of main.pdf) resolve from main.aux via the xr package, so this document
% must build AFTER main.pdf (the Makefile orders it). References in the
% other direction (main text -> this document) are prose "Section S<n>"
% pointers, never \ref — the labels live in separate PDFs.
%
% BUILD: `make supplement` (or plain `make`, which builds both PDFs).
% ===========================================================================
\documentclass[review,12pt,authoryear]{elsarticle}

\usepackage{amsmath}
\usepackage{amssymb}
\usepackage{booktabs}
\usepackage{longtable}
\usepackage{array}
\usepackage{url}
\usepackage{lineno}
\modulolinenumbers[5]
% Long unbreakable typewriter tokens (hashes, model IDs) need extra
% justification headroom (same rationale as main.tex).
\setlength{\emergencystretch}{3em}
\journal{Simulation Modelling Practice and Theory}
% Unsubmitted: the first-page footer reads "Preprint" (same as main.tex).
\makeatletter
\def\ps@pprintTitle{%
  \let\@oddhead\@empty
  \let\@evenhead\@empty
  \def\@oddfoot{\footnotesize\itshape Preprint\hfill\today}%
  \let\@evenfoot\@oddfoot}
\makeatother

% --------------------------------------------------------------------------
% Shared manuscript macros (subset of main.tex; \figstub is main-only —
% all figure artwork lives in the main document). Defined BEFORE
% \externaldocument: main.aux label entries carry caption text that uses
% these macros, and the xr import trips on undefined control sequences.
% --------------------------------------------------------------------------
\newcommand{\expid}[1]{\textsf{#1}}
\newcommand{\code}[1]{\texttt{\detokenize{#1}}}

% Cross-references into the main text resolve from main.aux. main.aux is
% written by a hyperref build (5-field \newlabel entries), so plain xr
% leaks the extra fields into the typeset text; xr-hyper + hyperref parse
% them correctly. xr-hyper must load BEFORE hyperref.
\usepackage{xr-hyper}
\usepackage{xcolor}
\definecolor{obLink}{RGB}{28,60,138}
\usepackage[colorlinks=true, allcolors=obLink, bookmarksnumbered=true,
            bookmarksopen=true, bookmarksopenlevel=1]{hyperref}
\externaldocument{main}

% Supplementary numbering: S-prefixed counters per the Elsevier
% online-supplement convention (sections S1..S5, tables S1.., equations S1..).
\renewcommand{\thesection}{S\arabic{section}}
\renewcommand{\thetable}{S\arabic{table}}
\renewcommand{\thefigure}{S\arabic{figure}}
\renewcommand{\theequation}{S\arabic{equation}}

\begin{document}

\begin{frontmatter}

\title{Supplementary Material for ``OpenBilliards: a validated, bit-exact
event-driven simulator for cue-sports mechanics, with a systematic
assessment of the empirical laws of pool''}

\author{OpenBilliards contributors}

\begin{abstract}
This document is the online supplement to the main article. Section~S1
maps every figure and table of the paper to its experiment, pinned
output hash, and registered source. Section~S2 states the tier-1
reproduction mechanics of the artifact bundle. Section~S3 reproduces the
governing equations and frozen constants of every model-catalog row.
Section~S4 carries the extended per-experiment result prose for
experiment groups B (stick--ball) and D (curved paths), whose summaries
appear in the main text. Section~S5 transcribes the full predicted-shift
verification table for the slip-law and branch-selection fix cascade.
Section~S6 documents the presentation methods behind the artifact
bundle's validated-scenario video pack --- the $C^2$ log-domain time
remap and shutter synthesis --- which render already-verified kernel
records and carry no validation claims.
Numbered cross-references to sections, tables, and figures without an
``S'' prefix resolve to the main article.
\end{abstract}

\end{frontmatter}

\linenumbers

% ---------------------------------------------------------------------------
% ===========================================================================
% Supplement S1 (former Appendix A) — Data provenance
% Compiles into supplement.pdf via supplement-main.tex; main-text pointers
% are prose "Section S1".
% Maps every figure/table to its experiment ID, baseline-manifest row
% (expected sha256, 8-char prefix shown; full hashes are in the
% machine-readable expected-hash table of the artifact bundle), and
% registered source. Slot namespace: figure slots F1-F14 vs experiment
% IDs F1-F3.
% ===========================================================================
\section{Data provenance: figures and tables to pinned outputs}
\label{sec:appendix-provenance}

Table~\ref{tab:provenance} states the paper's reading contract: every
plotted number regenerates from one registered script against the pinned
kernel, and the expected sha256 of each output is recorded both in the
study's baseline manifest (the post-cascade baseline of record under the
canonical kernel \code{0.1.0+fb3425ad}) and in the artifact bundle's
machine-readable expected-hash table, which is parsed from that manifest
rather than hand-written. The scripts, the manifest, and the bundle are
not part of the public source release (main-text
Section~\ref{sec:availability}). Slot F1 (the architecture schematic)
carries no experiment data. Tables~\ref{tab:models}, \ref{tab:domains}, and
\ref{tab:determinism} draw on battery output metadata, the group
subsections' ledgered numbers, and the determinism evidence index,
respectively.

{\footnotesize
\begin{longtable}{lllp{0.5\linewidth}}
  \caption{Provenance map: figure slot $\to$ experiment $\to$ expected
  output hash (prefix) $\to$ registered source.}
  \label{tab:provenance}\\
  \toprule
  Slot & Exp & sha256 (8) & Registered source \\
  \midrule
  \endfirsthead
  \toprule
  Slot & Exp & sha256 (8) & Registered source \\
  \midrule
  \endhead
  \bottomrule
  \endfoot
  F3  & \expid{A1} & \code{ec64b364} & TP 3.1, A.5, A.7 \\
  F3  & \expid{A2} & \code{865815ab} & TP 3.3, A.4 \\
  F2  & \expid{A3} & \code{f1367bc1} & TP A-14, A.28, B.3 \\
  F4  & \expid{A4} & \code{dcbe8050} & TP B.25, FAQ speed-effects \\
  F4  & \expid{A5} & \code{df91770b} & TP A.26 \\
  F4  & \expid{A6} & \code{c24e4e37} & TP A.27 \\
  F5  & \expid{A7} & \code{87bd5f07} & TP A.17, A.14 bottom \\
  F6  & \expid{B1} & \code{faebd431} & TP A-31; \citet{cross2008squirt} \\
  F6  & \expid{B2} & \code{3b6356e4} & TP B.1, B.26 \\
  F6  & \expid{B3} & \code{1ed54d97} & TP 2.1, A.12 \\
  F6  & \expid{B4} & \code{a42720a6} & TP 4.2 \\
  F7  & \expid{C1} & \code{0b3d1246} & TP B.28, B.27 \\
  F8  & \expid{C2} & \code{64410241} & TP 7.2; \citet{cohen2002threecushion} \\
  F8  & \expid{C3} & \code{3d12bdf4} & FAQ; TP 6.3, 7.3 \\
  F9  & \expid{C4} & \code{32feed7a} & \citet{mathavan2010cushion} Figs 7--10; HSV B.15 \\
  F10 & \expid{D1} & \code{8a4a53a2} & TP A.4; HSV A.76; \citet{coriolis1835} \\
  F10 & \expid{D2} & \code{2178b9d8} & TP A.19; \citet{coriolis1835} \\
  F10 & \expid{D3} & \code{c5a5113f} & TP B.23 \\
  F11--F12 & \expid{E1} & \code{21f9efd2} & Jewett via BD Dec 2024; HSV A.88--A.95 \\
  F11 & \expid{E2} & \code{ea6cff60} & TP B.21, A.15 \\
  F12 & \expid{E3} & \code{87121a5b} & HSV 7.6/7.7, B.43, B.45 \\
  F13 & \expid{F1} & \code{28ee6348} & TP 4.1, B.5, B.6, B.8 \\
  F13 & \expid{F2} & \code{78d1b730} & TP B.8/B.9; FAQ trisect \\
  F13 & \expid{F3} & \code{f0ee2200} & FAQ jump; HSV A.1/A.2, B.37 \\
  F14 & ORACLE-X   & \code{00a1488e} & \citet{kiefl2024pooltool}, pinned oracle env \\
\end{longtable}
}

The abstract's battery-scale figure --- 16{,}933 recorded simulations
--- is an exact, artifact-backed count, not a program-management
estimate: it is the number of kernel simulations the battery executes
end-to-end, recorded one row per call (serde-exact inputs plus
stamp-normalized trajectory digest) by the input tap behind the
wasm-replay corpus of Section~\ref{sec:reproducibility}, whose SHA-256
and per-experiment row counts are pinned in the bundle's replay report.
It supersedes the pre-replay floor of ``more than 4{,}800'', which was
derived from the pinned outputs alone: the baseline manifest's digest
normalization counts 3{,}157 digest-pinned per-record trajectory
fingerprints across the 25 outputs --- all 3{,}157 of which appear in the
replay corpus's digest set --- and the remainder are digest-free recorded sweep rows (e.g.\ the
492-run \expid{A4} grid) carried in the same pinned files.

One artifact is deliberately outside this contract. The computational
performance evidence of Section~\ref{sec:simulator-performance}
(Table~\ref{tab:perf}) is engineering-tier, not physics V\&V: its
timings are machine-dependent by nature, so its artifact (a pinned
benchmark report in the bundle) records host provenance, both
measurement passes, and the protocol, but
carries \emph{no} expected hash and is excluded from the physics-identity
manifest. Its determinism guards are digest ties, not timing hashes: the
timed kernel records reproduce the stamp-independent trajectory digests
of the ORACLE-X output (30/30 matched), proving the timed scenarios are
the event-verified ones.

Three run-order dependencies exist; every other script is order-free.
Experiment \expid{A7} reads the \expid{A3} output; ORACLE-X reads the
\expid{A3} and \expid{A6} outputs and additionally requires the pinned
pooltool environment; \expid{D1} and \expid{F2} re-read only their own
outputs. Digitized-source fixtures (the TP curves, the Mathavan figures,
the FAQ charts) carry per-fixture provenance blocks and digitization
uncertainty; no third-party source frames are redistributed.

% ===========================================================================
% Supplement S2 (former Appendix B) — Tier-1 reproduction procedure
% Compiles into supplement.pdf via supplement-main.tex; main-text pointers
% are prose "Section S2".
% Cites the bundle manifest + checker; step-by-step tier-1 procedure;
% what "verified" means mechanically. Sources: the bundle's results map
% (contents, how to reproduce, caveats), the study design's tier-1
% definition, the ledger's exclusion sections, the identity checker, and
% the determinism matrix.
% ===========================================================================
\section{Tier-1 reproduction procedure}
\label{sec:appendix-repro}

\paragraph{Artifact bundle contents.}
The artifact bundle, which is not part of the public source release and
has not been published (main-text Section~\ref{sec:availability}),
contains: the canonical WebAssembly binary (\code{pool_wasm_bg.wasm},
stamp \code{0.1.0+fb3425ad}, sha256
\code{5e3b6ee499b32c868689c3f5279bbd02}\allowbreak\code{8161c94bdd306fbe969a2dbe5ef34936});
a binary Python wheel (abi3, CPython $\geq$ 3.9) and an npm SDK tarball;
the 24 experiment scripts, the oracle harness, the shared experiment
library, and a 13-check environment smoke test; the sensitivity pack
(the sensitivity scripts, the report generator, and the pinned
sensitivity report) together with the physics-identity
checker it imports and the checker's normalized-baseline manifest and
machine waiver ledger; exactly the fixtures the
scripts read, each with its provenance block; the machine-readable
expected-hash table, parsed from the baseline manifest and never
hand-written; the waiver ledger verbatim; the validated-scenario video
pack --- sixteen annotated
slow-motion entries (final renders 2026-08-05), each staged strike
replayed through the then-current shipped canonical binary under a
gate-before-capture law (the claimed pinned record digest must
reproduce before any frame is captured; the per-entry kernel stamp is
recorded in the gate report, and the pack's stamps are later canonical
stamps than the battery pin above, each entry's digests pinned against
its own stamp), with a per-entry machine-readable gate report
(the derived frame-schedule summary --- frame
counts, worst rate-step ratio, shutter parameters, replay spans ---
gate verdicts, and, for part
of the pack, capture-verification records; the playback plans
themselves are authored in the render tool's scenario catalog), a
poster frame, and a pack
ledger carrying per-entry SHA-256 hashes of the delivered bytes
together with the render and verification history, including the
cue-model provenance (presentation methods
in Section~\ref{sec:supp-timeremap}; the videos present verified
records and carry no claims of the paper) --- and
a SHA-256 manifest covering every file in the bundle.

\paragraph{Procedure.}
Tier-1 reproduction is five steps: (1)~verify the SHA-256 manifest;
(2)~install the wheel; (3)~run the 13-check environment smoke test;
(4)~run any experiment script; (5)~compare the output sha256 against the
expected-hash table. Verification \emph{is} hash equality: a tier-1
reproduction succeeds exactly when the regenerated output bytes hash to
the pinned value.

\paragraph{Checker semantics.}
The physics-identity checker (normalization spec v2)
masks only kernel-version stamps and compares per-record trajectory
digests as stamp-independent physics evidence; 25/25 outputs pass on the
baseline manifest of record.

\paragraph{Reproducibility caveats, stated honestly.}
Wheel bytes are not guaranteed bit-reproducible across machines and
toolchains --- the determinism obligation attaches to \emph{outputs}, not
to package bytes. Raw output hashes embed the kernel stamp, so a
differently-stamped wheel requires the documented normalization until the
submission-freeze rebuild.

\paragraph{Licensing.}
The kernel source is public under the MIT license, and the portions
derived from pooltool remain under the Apache License~2.0 (main-text
Section~\ref{sec:availability}); the bundle's license terms will be
stated when it is published. The bundle redistributes no Dr.~Dave
frames or other third-party captures --- only derived measurement
tables with provenance.

% ===========================================================================
% Supplement S3 (former Appendix C) — Governing equations and frozen constants
% Compiles into supplement.pdf via supplement-main.tex; main-text pointers
% are prose "Section S3".
% The physics must be recoverable from the paper alone, independent of any
% source revision. This appendix states every validated model's equations
% as implemented, with published citations and itemized deviations, and
% consolidates every frozen constant with its provenance. No source code
% appears here; the equations below are transcriptions of the canonical
% binary's arithmetic, not re-derivations.
% ===========================================================================
\section{Governing equations and frozen constants of the validated models}
\label{sec:appendix-equations}

This appendix makes each row of Table~\ref{tab:models} reproducible from
the paper alone. Every equation is stated \emph{as implemented at kernel
\code{0.1.0+fb3425ad}} (the canonical binary of
Section~\ref{sec:simulator-determinism}); the citation attached to each
model identifies its published lineage, and wherever the implementation
deviates from the cited source the deviation is itemized in place rather
than left behind the word ``corrected.'' Formulations inherited from the
open-source reference implementation pooltool \citep{kiefl2024pooltool}
are cited as such; the oracle comparison of Section~\ref{sec:oracle}
independently verifies that inheritance event-by-event \emph{on the
matched registry only}---the 38 registered \expid{A3}/\expid{A6}
points, which exercise the Han cushion family, the legacy ball--ball
law, and matched strike models. The Stronge, Mathavan, and
compliant-cluster rows are never oracle-verified (Section~\ref{sec:oracle}
states this explicitly for Stronge); their evidence is the
published-curve anchors and battery gates of
Sections~\ref{sec:results} and~\ref{sec:oracle}. All frozen
constants appear in Table~\ref{tab:frozen-constants} with units and
provenance.

Notation: $\mathbf{v}$ and $\boldsymbol{\omega}$ are a ball's
center-of-mass velocity and angular velocity, $R$ its radius, $m$ its
mass, $I = \tfrac{2}{5}mR^{2}$ its moment of inertia,
$\hat{\mathbf{z}}$ the vertical unit vector, and $g = 9.81$\,m/s$^2$.
Subscripts $n$ and $t$ denote components along a contact normal
$\hat{\mathbf{n}}$ and tangent $\hat{\mathbf{t}}$. A hat denotes a unit
vector; a prime denotes a post-impact quantity.

\subsection{Ball--cloth motion between events}
\label{sec:appendix-eq-cloth}

As implemented at kernel \code{0.1.0+fb3425ad}, cloth contact uses three
independent friction regimes---sliding $\mu_s$, rolling $\mu_r$, and
vertical-spin decay $\mu_{sp}$---with closed-form kinematics in each
regime, following the classical treatment collected in the
technical-proof corpus \citep{alciatore_tp_catalog} and the event-based
formulation of \citet{leckie2006event} as realized by
\citet{kiefl2024pooltool}. The cloth contact-point slip velocity is
\begin{equation}
\mathbf{u} \;=\; \mathbf{v} + R\,(\hat{\mathbf{z}} \times
\boldsymbol{\omega}).
\label{eq:cloth-slip}
\end{equation}

\paragraph{Sliding} While $\mathbf{u} \neq \mathbf{0}$, planar Coulomb
friction acts along the (exactly constant) initial slip direction
$\hat{\mathbf{u}}_0$:
\begin{align}
\mathbf{r}(t) &= \mathbf{r}_0 + \mathbf{v}_0 t
  - \tfrac{1}{2}\,\mu_s g\,t^{2}\,\hat{\mathbf{u}}_0, &
\mathbf{v}(t) &= \mathbf{v}_0 - \mu_s g\,t\,\hat{\mathbf{u}}_0,
\label{eq:slide-rv}\\
\boldsymbol{\omega}_{xy}(t) &= \boldsymbol{\omega}_{0,xy}
  - \frac{5\,\mu_s g}{2R}\,t\,(\hat{\mathbf{u}}_0 \times
  \hat{\mathbf{z}}), &
\mathbf{u}(t) &= \mathbf{u}_0 - \tfrac{7}{2}\,\mu_s g\,t\,
  \hat{\mathbf{u}}_0.
\label{eq:slide-w}
\end{align}
The constancy of $\hat{\mathbf{u}}_0$ is an exact consequence of planar
Coulomb friction, not an approximation; the slide ends when
$\mathbf{u} = \mathbf{0}$, at
\begin{equation}
\tau_{\mathrm{slide}} = \frac{2\,|\mathbf{u}_0|}{7\,\mu_s g}.
\label{eq:slide-tau}
\end{equation}

\paragraph{Rolling} A rolling ball travels in a straight line along
$\hat{\mathbf{v}}_0$ with constant deceleration $\mu_r g$, and the
rolling constraint is re-imposed exactly at every instant (not
integrated):
\begin{equation}
\mathbf{v}(t) = \mathbf{v}_0 - \mu_r g\,t\,\hat{\mathbf{v}}_0,
\qquad
\boldsymbol{\omega}_{xy}(t) \;=\;
\frac{\hat{\mathbf{z}} \times \mathbf{v}(t)}{R},
\qquad
\tau_{\mathrm{roll}} = \frac{|\mathbf{v}_0|}{\mu_r g}.
\label{eq:roll}
\end{equation}

\paragraph{Vertical spin} $\omega_z$ decays linearly and independently
in every \emph{cloth-contact} moving state (sliding, rolling,
spinning), clamped at exactly zero (no overshoot)---airborne flight is
excluded: it freezes all of $\boldsymbol{\omega}$ (below), so a
ball in flight lands with the spin it launched with:
\begin{equation}
\omega_z(t) = \operatorname{sgn}(\omega_{z,0})\,
\max\!\bigl(|\omega_{z,0}| - \alpha t,\; 0\bigr),
\qquad
\alpha = \frac{5\,\mu_{sp}\,g}{2R},
\label{eq:spin-decay}
\end{equation}
with the frozen $\mu_{sp} = \tfrac{4}{9}R$ giving
$\alpha = \tfrac{10}{9}g \approx 10.9$\,rad/s$^2$. A units convention
governs this line: $\mu_{sp}$ is \emph{dimensionless} --- the
inherited shorthand $\tfrac{4}{9}R$ consumes the numeric value of $R$
in metres, i.e.\ $\mu_{sp} = \tfrac{4}{9}\,(R/\mathrm{1\,m}) = 0.0127$
--- so $\alpha = 5\mu_{sp}g/(2R)$ carries s$^{-2}$ as required; both
the coefficient and the resulting rate are printed so the
implementation is reproducible either way. The spin regime ends
at $\tau_{\mathrm{spin}} = |\omega_z|/\alpha$; a ball is projected to
exact rest (all components identically zero) when its last regime
expires.

\paragraph{Airborne flight} A ball off the bed follows an exact
ballistic parabola, $\mathbf{r}(t) = \mathbf{r}_0 + \mathbf{v}_0 t -
\tfrac{1}{2} g t^{2}\hat{\mathbf{z}}$, with $\boldsymbol{\omega}$
frozen: the design excludes air drag and Magnus lift so that every
trajectory stays polynomial (Section~\ref{sec:simulator-core}). Every
landing resolves through the slate impact of
\ref{sec:appendix-eq-slate}.

\subsection{Ball--slate impact (landing resolution)}
\label{sec:appendix-eq-slate}

As implemented at kernel \code{0.1.0+fb3425ad}, every landing of an
airborne ball on the bed---including the immediate post-strike bounce
that every jump shot dispatches to at $\tau \approx 0$
(\ref{sec:appendix-eq-strike})---resolves through a frictional
inelastic impulse in the formulation of \citet{kiefl2024pooltool}. The
ball is first re-seated to exactly $z = R$. With $v_z < 0$ the incoming
vertical velocity, the normal channel is restitution $e_t$:
\begin{equation}
v'_z \;=\; -\,e_t\,v_z
\qquad\text{(normal impulse } (1 + e_t)\,|v_z| \text{ per unit mass)}.
\label{eq:slate-normal}
\end{equation}
The tangential channel is resolved with the normal problem settled
first: the cloth contact slip $\mathbf{u} = \mathbf{v}_{xy} +
R\,(\hat{\mathbf{z}} \times \boldsymbol{\omega})$
(Eq.~\eqref{eq:cloth-slip}, planar velocity only) selects between the
Coulomb slip impulse and the $2/7$-stick impulse, stick applying when
there is no slip or the stick impulse is not larger in magnitude:
\begin{equation}
\Delta\mathbf{V}_t =
\begin{cases}
-\,\mu_s\,(1+e_t)\,|v_z|\;\hat{\mathbf{u}}, &
\mathbf{u} \neq \mathbf{0} \text{ and }
|\Delta\mathbf{V}_{t,\mathrm{stick}}| >
|\Delta\mathbf{V}_{t,\mathrm{slip}}|
\quad\text{(slip)},\\[4pt]
\tfrac{2}{7}\bigl[R\,(\boldsymbol{\omega} \times \hat{\mathbf{z}}) -
\mathbf{v}_{xy}\bigr], & \text{otherwise (stick)},
\end{cases}
\label{eq:slate-tangential}
\end{equation}
with the corresponding spin updates
\begin{equation}
\text{slip:}\quad
\boldsymbol{\omega}' = \boldsymbol{\omega} +
\frac{5}{2R}\,|\Delta\mathbf{V}_t|\,
(\hat{\mathbf{z}} \times \hat{\mathbf{u}}),
\qquad\quad
\text{stick:}\quad
\boldsymbol{\omega}' = \tfrac{2}{7}\,\boldsymbol{\omega} +
\tfrac{5}{7}\,\frac{\hat{\mathbf{z}} \times \mathbf{v}_{xy}}{R}
\label{eq:slate-spin}
\end{equation}
(note the stick update also scales $\omega_z$ by $2/7$). That
$\omega_z$ scaling deserves an explicit acknowledgment rather than a
parenthesis: a tangential impulse applied at the contact point directly
below the ball centre exerts no torque about the vertical axis, so the
$5/7$ vertical-spin loss per stuck bounce is not derivable from the
printed impulse --- it is an inherited artifact of the ported
formulation \citep{kiefl2024pooltool}, disclosed here as a known
limitation: every jump shot dispatches an immediate slate bounce, so a
landing that resolves through the stick branch bleeds side spin
accordingly. The friction
coefficient is the cloth sliding $\mu_s$---the bounce happens through
the cloth. Finally the bounce guard: if the outgoing apex
$v'^{2}_z/2g$ falls below the keep threshold---$5$\,mm under the
parity default, $10\,\mu$m under the three realism presets, the same
two-regime convention as the ball--ball micro-hop guard of
\ref{sec:appendix-eq-ballball}---then $v'_z := 0$ and the bounce
cascade terminates (the apex decays geometrically with $e_t^{2}$). The
frozen $e_t = 0.5$ is the reference default
(Table~\ref{tab:frozen-constants}); no published measurement exists.

\subsection{Ball--ball impact}
\label{sec:appendix-eq-ballball}

As implemented at kernel \code{0.1.0+fb3425ad} (catalog row
\code{CoulombCap}), the ball--ball impact is an instantaneous impulse
with a normal restitution channel and a tangential friction channel,
following the throw model of the technical-proof corpus
\citep{alciatore_tp_a14,alciatore_tp_catalog} in the formulation of
\citet{kiefl2024pooltool}. The contact normal $\hat{\mathbf{n}}$ is the
full three-dimensional line of centers from ball~1 to ball~2.

\paragraph{Normal channel} With $e_b$ the arithmetic mean of the two
balls' restitution parameters ($e_b = 0.95$ for the catalog ball set),
the normal velocity components exchange as an equal-mass
one-dimensional inelastic collision:
\begin{equation}
v'_{1n} = \tfrac{1}{2}\bigl[(1-e_b)\,v_{1n} + (1+e_b)\,v_{2n}\bigr],
\qquad
v'_{2n} = \tfrac{1}{2}\bigl[(1+e_b)\,v_{1n} + (1-e_b)\,v_{2n}\bigr].
\label{eq:bb-normal}
\end{equation}

\paragraph{Friction coefficient} The tangential channel samples the
velocity-dependent friction curve of \citet{alciatore_tp_a14} at the
contact-point relative surface speed
$u_{12} = |\mathbf{u}_{12}|$, where $\mathbf{u}_{12}$ is the relative
velocity of the two contacting surface points (both balls' spins
included):
\begin{equation}
\mu_{bb}(u_{12}) \;=\; 0.009951 + 0.108\,e^{-1.088\,u_{12}}.
\label{eq:bb-mu}
\end{equation}
(The revision drift within the source corpus for the decay constant is
ledgered as a source discrepancy, Section~\ref{sec:simulator-catalog}.)

\paragraph{Slip branch (Coulomb cap)} If the contact is slipping, each
ball receives an equal-and-opposite tangential velocity change of
magnitude $\mu_{bb}\,\hat{J}_n$ opposing the slip, with the
per-unit-mass normal impulse
\begin{equation}
\hat{J}_n \;=\; \frac{1+e_b}{2}\,\bigl|v_{1n} - v_{2n}\bigr|
\qquad \text{(TP~A.5, Eq.~10; catalog law \code{CoulombCap})},
\label{eq:bb-cap}
\end{equation}
\begin{align}
\Delta\mathbf{V}_{t,1} &= -\,\mu_{bb}(u_{12})\,\hat{J}_n\,
\hat{\mathbf{u}}_{12},
\qquad
\Delta\mathbf{V}_{t,2} = -\Delta\mathbf{V}_{t,1},
\label{eq:bb-slip}\\
\Delta\boldsymbol{\omega}_{1} &= \Delta\boldsymbol{\omega}_{2}
= \frac{5}{2R}\,\hat{\mathbf{n}} \times \Delta\mathbf{V}_{t,1}.
\nonumber
\end{align}
In a stun cut this yields a slip-regime throw angle of exactly
$\arctan \mu_{bb}$, independent of $e_b$. The \emph{legacy} law
(recorded as \code{PostImpactNormalDelta}; retained solely for
byte-exact replay of pre-correction records and oracle parity,
Section~\ref{sec:oracle}) scales the same impulse with
$e_b\,|v_{1n}-v_{2n}|$ instead, leaving slip-regime throw
$\arctan[2\mu_{bb}e_b/(1+e_b)]$---about 2.6\% low at $e_b=0.95$ (the
slip-law defect).

\paragraph{Stick branch and boundary} If the contact has no initial
slip, or the slip solution of Eq.~\eqref{eq:bb-slip} would reverse the
slip direction ($\mathbf{u}_{12}\cdot\mathbf{u}'_{12} \le 0$), the
no-slip (``$1/7$'') solution applies instead:
\begin{align}
\Delta\mathbf{V}_{t,1} &= -\tfrac{1}{7}\Bigl[(\mathbf{v}_1 -
\mathbf{v}_2)_t + R\,(\boldsymbol{\omega}_1 +
\boldsymbol{\omega}_2)\times\hat{\mathbf{n}}\Bigr],
\label{eq:bb-stick}\\
\Delta\boldsymbol{\omega}_{1,2} &=
-\tfrac{5}{14}\Bigl[\frac{\hat{\mathbf{n}}\times(\mathbf{v}_1 -
\mathbf{v}_2)}{R} + (\boldsymbol{\omega}_1 +
\boldsymbol{\omega}_2)_t\Bigr].
\nonumber
\end{align}
The stick/slip boundary is therefore implicit,
$u_{12} = 7\,\mu_{bb}(u_{12})\,\hat{J}_n$, and moves with the cap: the
corrected law's boundary cut sits $0.6$--$1^\circ$ above the legacy
law's (the predicted-shift verification of Section~\ref{sec:oracle}).
Spin components along $\hat{\mathbf{n}}$ are unchanged by the impact.

\paragraph{Vertical micro-hop guard} The three-dimensional line of
centers gives the friction impulse a vertical component (ball hop). An
on-bed ball keeps its post-impact $v_z$ only when the resulting bounce
apex $v_z^{2}/2g$ clears a threshold: $5$\,mm under the parity default
(absorbing the micro-hop, matching the planar reference formulation)
and $10\,\mu$m under the three realism presets of
Table~\ref{tab:models} (\code{stronge_realistic},
\code{mathavan_fitted}, \code{grip_fitted}), which keep the physical
micro-hop while absorbing numerical dust.

\subsection{Stick--ball strike}
\label{sec:appendix-eq-strike}

As implemented at kernel \code{0.1.0+fb3425ad}, the strike is the
instantaneous-point cue impact of the technical-proof corpus (TP~A-30)
with the endmass squirt model (TP~A-31)
\citep{alciatore_tp_a31,cross2008squirt,shepard2001squirt}, in the
formulation of \citet{kiefl2024pooltool}. Shot parameters are the cue
speed $v_0 \in (0, 15]$\,m/s, azimuth $\varphi$, elevation
$\theta \in [0, \pi/2)$, and the commanded \emph{tip-center} offsets
$(a_{\mathrm{cmd}}, b_{\mathrm{cmd}})$ in units of $R$ (side, vertical), subject to the miscue envelope
\begin{equation}
a_{\mathrm{cmd}}^{2} + b_{\mathrm{cmd}}^{2} \;\le\; 0.25
\qquad\text{(offset radius $\le R/2$)};
\label{eq:strike-miscue}
\end{equation}
shots beyond the envelope are rejected as invalid input, never
resolved. Post-registration, a kernel amendment (kernel
\code{0.1.0+8efd660e}) re-binds this limit to the \emph{contact}
offset --- the variable in which the $0.5$--$0.55R$ literature bracket
is stated --- raising the commanded cap to $0.5\,(1 + \rho/R)$
($0.686R$, nickel tip); every previously accepted input reproduces
bit-exactly (\expid{B3}'s registered nickel-tip rejection probes,
which would now be accepted, are the one disclosed behavioral change
class), so Eq.~\eqref{eq:strike-miscue} is the envelope of record for
every result in this paper (amendment note: main-text
Section~\ref{sec:limitations}).

\paragraph{Commanded-to-internal offset map (frame and sign)} Two
side-offset conventions exist and must not be conflated. The
\emph{commanded} offset $a_{\mathrm{cmd}}$ of the shot record is
positive for \emph{right} English (the shooter's perspective); the
\emph{internal} offset $a$ consumed by every formula below is measured
toward the shooter's \emph{left} ($+\hat{\mathbf{x}}$ in the
cue-aligned frame). The resolver's first act is therefore a sign
inversion, fused with the tip-curvature contact correction (tip-center
offsets map to the contact point on the ball surface through the tip
curvature radius $\rho$):
\begin{equation}
a \;=\; \frac{-\,a_{\mathrm{cmd}}}{1 + \rho/R},
\qquad
b \;=\; \frac{b_{\mathrm{cmd}}}{1 + \rho/R},
\label{eq:strike-tipscale}
\end{equation}
with $\rho = 0.0106$\,m (nickel) for the default cue, and
$c = \sqrt{1 - a^{2} - b^{2}}$ is the contact-normal component
(offsets internal and tip-corrected hereafter). A reproduction that
scales but does not negate the side offset mirrors both the spin of
Eq.~\eqref{eq:strike-spin} and the squirt direction of
Eq.~\eqref{eq:strike-squirt} on every side-English shot.

\paragraph{Ball speed} With cue mass $M$, ball mass $m$, and tip--ball
restitution $e_{\mathrm{tip}}$, the post-strike ball speed is
\begin{equation}
v \;=\; \frac{(1 + e_{\mathrm{tip}})\,v_0}
{1 + \dfrac{m}{M} + \dfrac{5}{2}\,(a^{2} + b^{2})},
\label{eq:strike-speed}
\end{equation}
the elevation-independent (ball-frame) form of TP~A-30; the frozen
default $e_{\mathrm{tip}} = 1$ reproduces the classical $2v_0$
numerator, and the cue catalog carries measured tip values (leather
$0.73$, phenolic $0.85$; Table~\ref{tab:frozen-constants}). In the
cue-aligned frame (cue axis along $-\hat{\mathbf{y}}$, elevated by
$\theta$ about $\hat{\mathbf{x}}$; $a$ measured toward
$+\hat{\mathbf{x}}$, the shooter's left; $b$ upward) the outgoing
velocity is $(0, -v\cos\theta, -v\sin\theta)$ and the angular velocity
is
\begin{equation}
\boldsymbol{\omega}_{\mathrm{cue}} \;=\; \frac{5v}{2R}\,
\bigl(b,\; a\sin\theta,\; -a\cos\theta\bigr),
\label{eq:strike-spin}
\end{equation}
both rotated to the table frame by $\varphi + \pi/2$ about
$\hat{\mathbf{z}}$. Elevation drives the ball downward
($v_z = -v\sin\theta$); for $0 < \theta < \theta_{\mathrm{full}}$ the
ball leaves the strike airborne---the pinch of
Eq.~\eqref{eq:strike-pinch} removes only the fraction $\lambda < 1$ of
the downward component---and its immediate slate bounce
(\ref{sec:appendix-eq-slate}) is the legal-jump mechanism. At full
pinch, $\theta \ge \theta_{\mathrm{full}}$, $\lambda = 1$ and the bed
reaction removes the \emph{entire} downward component before
motion-state classification: the strike exits grounded, with no slate
bounce.

\paragraph{Squirt (cue-ball deflection)} The endmass model of TP~A-31
(eq.~13 there; cf.\ \citealp{shepard2001squirt,cross2008squirt})
rotates the outgoing \emph{velocity} (spin untouched) away from the
aiming line by the squirt angle $\psi$, directed opposite the applied
side spin:
\begin{equation}
\tan\psi \;=\; \frac{\tfrac{5}{2}\,a\,\sqrt{1-a^{2}}}
{1 + \dfrac{m}{m_e} + \dfrac{5}{2}\,(1-a^{2})},
\label{eq:strike-squirt}
\end{equation}
with $m_e$ the effective endmass. The default cue freezes
$m_e = m/30 \approx 5.67$\,g; the per-shaft catalog spans
$1.70$--$12.7$\,g through the small-offset pivot-length inversion
$m/m_e = 2.5\,L_p/R - 3.5$ of \citet{shepard2001squirt}, with
$L_p$ the robot-measured pivot length
(Table~\ref{tab:frozen-constants}).

\paragraph{Pinch (elevated-strike) regime} Below
$\theta_{\mathrm{onset}} = 60^{\circ}$ the strike is exactly
Eqs.~\eqref{eq:strike-speed}--\eqref{eq:strike-squirt}. Above it the
tip pins the ball against the cloth (the grounded-mass\'e regime of
\citealp{coriolis1835}) and the bed reaction participates in the
strike, with engagement fraction
\begin{equation}
\lambda(\theta) = s^{2}(3 - 2s), \qquad
s = \operatorname{clamp}\!\Bigl(
\frac{\theta - \theta_{\mathrm{onset}}}
{\theta_{\mathrm{full}} - \theta_{\mathrm{onset}}},\, 0,\, 1\Bigr),
\qquad \theta_{\mathrm{full}} = 70^{\circ},
\label{eq:strike-pinch}
\end{equation}
a $C^{1}$ smoothstep. The bed absorbs the normal launch impulse
$\hat{\jmath} = \lambda(\theta)\,|v_z|$ per unit mass
($v_z \to (1-\lambda)v_z$), and the same impulse drags cloth friction
$\mu_s\,\hat{\jmath}$ through the contact patch opposite the
contact-point slip, with exactly the slip/stick selection of
Eq.~\eqref{eq:slate-tangential} (normal impulse magnitude
$\hat{\jmath}$ in place of $(1+e_t)|v_z|$). The pinch
leaves the ball pressed into the cloth: for a window
$\tau_{\mathrm{pinch}} = 0.12$\,s the sliding friction deceleration is
elevated to $\mu_s k g$ with
\begin{equation}
k \;=\; 1 + \frac{\kappa\,\hat{\jmath}}{g\,\tau_{\mathrm{pinch}}},
\qquad \kappa = 0.45,
\label{eq:strike-press}
\end{equation}
each phase remaining an exact parabola. The constants
$(\tau_{\mathrm{pinch}}, \kappa)$ are first-party fits to reference
mass\'e footage (Table~\ref{tab:frozen-constants}).

\subsection{Ball--cushion impact}
\label{sec:appendix-eq-cushion}

All four cushion rows of Table~\ref{tab:models} share the contact
geometry and the two velocity-response hooks; they differ in the
impulse model. Throughout, $h$ is the cushion-nose height above the
bed and the nose contact angle $\theta_a$ satisfies
\begin{equation}
\sin\theta_a \;=\; \frac{h}{R} - 1 .
\label{eq:cushion-geom}
\end{equation}
The US-pool (WPA) geometry freezes $h = 0.635 \cdot 2R = 1.27R$, so
$\sin\theta_a = 0.27$; the snooker-lineage geometry of
\citet{mathavan2010cushion} places the contact at $h = 7R/5$
($\sin\theta_a = 0.4$) and is paired exclusively with the snooker
table configuration (the domain separation of
Section~\ref{sec:simulator-catalog}).

\paragraph{Velocity-dependent restitution hooks} Two separate restitution
definitions exist and are never conflated. The \emph{kinematic}
restitution consumed by the Han and Stronge models is
\begin{equation}
e(v_n) \;=\; \operatorname{clamp}\bigl(e_0 + k_v\,(v_n -
v_{\mathrm{ref}}),\; e_{\min},\; e_{\max}\bigr),
\label{eq:ec-kinematic}
\end{equation}
with $v_n$ the normal impact speed; the frozen default is constant
($e_0 = 0.85$, $k_v = 0$, $v_{\mathrm{ref}} = 1$\,m/s, clamp
$[0.6, 0.98]$). The \emph{energetic} (Stronge-sense) restitution
consumed by the Mathavan model shares the same functional form and
clamps but pivots on its own base $ee_0$:
\begin{equation}
ee(v_n) \;=\; \operatorname{clamp}\bigl(ee_0 + k_v\,(v_n -
v_{\mathrm{ref}}),\; e_{\min},\; e_{\max}\bigr),
\label{eq:ec-energetic}
\end{equation}
with $ee_0 = 0.98$ at the structural default
\citep[the snooker fit of][]{mathavan2010cushion}, $ee_0 = 0.80$
constant for \code{mathavan_fitted}, and the velocity-sloped
\code{grip_fitted} response
\begin{equation}
ee(v_n) = \operatorname{clamp}\bigl(0.65 + 0.18\,(v_n - 1),\;
0.5,\; 0.98\bigr).
\label{eq:ec-grip}
\end{equation}

\paragraph{Slip-velocity friction hooks} Cushion-nose
friction is an affine function of the contact-point slip speed
$v_{\mathrm{slip}}$ at impact start, clamped:
\begin{equation}
\mu_w(v_{\mathrm{slip}}) \;=\; \operatorname{clamp}\bigl(\mu_0 +
k\,(v_{\mathrm{slip}} - v_{\mathrm{ref}}),\; \mu_{\min},\;
\mu_{\max}\bigr).
\label{eq:muw-hook}
\end{equation}
The default row \code{han2005_default} freezes the flat value
$\mu_w = 0.20$ ($k = 0$). The \code{stronge_realistic} row carries the
video-calibrated curve $\mu_w = \operatorname{clamp}(0.20 -
0.4\,(v_{\mathrm{slip}} - 1),\, 0.20,\, 0.40)$: flat $0.20$ for slip
$\ge 1$\,m/s, rising toward $0.40$ for slow, spin-dominated contacts.
Both fitted Mathavan rows carry the same slope with re-fitted clamps,
$\mu_w = \operatorname{clamp}(0.20 - 0.4\,(v_{\mathrm{slip}} - 1),\,
0.16,\, 0.28)$. No published measurement of pool cushion
$\mu_w(v_{\mathrm{slip}})$ exists; the slopes are declared first-party
calibration content (Section~\ref{sec:simulator-catalog}).

\subsubsection{Corrected Han impulse model (han2005\_default)}
\label{sec:appendix-eq-han}

As implemented at kernel \code{0.1.0+fb3425ad}, the default cushion
model is the rigid-body impulse model of \citet{han2005dynamics},
evaluated in a cushion frame with $\hat{\mathbf{x}}$ the horizontal
normal from ball toward cushion, $\hat{\mathbf{y}}$ along the rail,
and $\hat{\mathbf{z}}$ up. With $(v_x, v_y, v_z)$ and
$(\omega_x, \omega_y, \omega_z)$ the incoming components in that frame
($v_z \equiv 0$ for a bed-contact ball), the contact-point slip at the
nose is (Han eqs.~14)
\begin{equation}
s_x = v_x \sin\theta_a - v_z\cos\theta_a + R\,\omega_y, \qquad
s_y = -v_y - R\,\omega_z \cos\theta_a + R\,\omega_x \sin\theta_a,
\label{eq:han-slip}
\end{equation}
with slip speed $s_0 = \sqrt{s_x^2 + s_y^2}$, and the normal-channel
scalar is $c = -v_x \cos\theta_a$. With the paper's coefficients
$A = 7/2m$ and $B = 1/m$ (Han eqs.~16), the normal impulse and the
slip-arrest impulse are (Han eqs.~17 and 20)
\begin{equation}
P_z^{E} = -\frac{(1+e)\,c}{B} = (1+e)\,m\,v_x \cos\theta_a,
\qquad
P_z^{S} = \frac{s_0}{A},
\label{eq:han-impulses}
\end{equation}
with $e = e(v_x)$ from Eq.~\eqref{eq:ec-kinematic} and
$\mu = \mu_w(s_0)$ from Eq.~\eqref{eq:muw-hook}. The case selection
and tangential impulses are
\begin{equation}
(P_x^{E}, P_y^{E}) =
\begin{cases}
\bigl(s_x/A,\; s_y/A\bigr), &
P_z^{S} \le \mu\,P_z^{E} \quad\text{(Han eqs.~18),}\\[4pt]
\bigl(\mu P_z^{E}\,s_x/s_0,\; \mu P_z^{E}\,s_y/s_0\bigr), &
\text{otherwise (Han eqs.~19),}
\end{cases}
\label{eq:han-cases}
\end{equation}
(the first case is sliding-then-adhesion, the second gross slip
throughout the impact), transformed to the rail frame (Han eqs.~21--22) as
$P_X = -P_x^{E}\sin\theta_a - P_z^{E}\cos\theta_a$,
$P_Y = P_y^{E}$,
$P_Z = P_x^{E}\cos\theta_a - P_z^{E}\sin\theta_a$, and applied as
(Han eqs.~23)
\begin{equation}
\begin{aligned}
v'_x &= v_x + P_X/m, &\quad
\omega'_x &= \omega_x - \tfrac{R}{I} P_Y \sin\theta_a,\\
v'_y &= v_y + P_Y/m, &
\omega'_y &= \omega_y + \tfrac{R}{I}\,(P_X \sin\theta_a -
P_Z\cos\theta_a),\\
v'_z &= v_z, &
\omega'_z &= \omega_z + \tfrac{R}{I} P_Y \cos\theta_a .
\end{aligned}
\label{eq:han-update}
\end{equation}
Pocket-jaw arcs use identical physics with the normal taken from the
arc center to the ball center.

\emph{Itemized deviations from the published paper} (the content of
``corrected'' in Table~\ref{tab:models}):
\begin{enumerate}
\item \textbf{Slip/stick case-selection inequality.} As printed, the
  paper's boundary between the final-slip and final-stick cases omits
  the friction coefficient; the implemented test is
  $P_z^{S} \le \mu\,P_z^{E}$ as in Eq.~\eqref{eq:han-cases}. Without
  the $\mu$, near-normal impacts select the wrong branch. The
  correction follows the reference implementation
  \citep{kiefl2024pooltool}, whose revision history documents it.
\item \textbf{Slip-direction angle $\theta_0$.} The paper's eqs.~18--19
  are written with $\sin\theta_0$, $\cos\theta_0$ for an angle that is
  never defined. The implementation fixes it to the initial
  contact-point slip direction, $\cos\theta_0 = s_x/s_0$,
  $\sin\theta_0 = s_y/s_0$, which is how it enters
  Eq.~\eqref{eq:han-cases} \citep[same lineage;][]{kiefl2024pooltool}.
\item \textbf{Restitution and friction inputs.} The paper's own
  speed-dependent restitution fit and incidence-dependent friction are
  not used; restitution and friction enter through the declared hooks
  of Eqs.~\eqref{eq:ec-kinematic} and \eqref{eq:muw-hook} (constant
  $e = 0.85$, $\mu_w = 0.20$ in the frozen default row).
\item \textbf{Planar output.} The vertical impulse $P_Z$ is computed
  (it feeds the $\omega_y$ update) but not applied to $v_z$; the
  vertical kinetic energy it would impart is treated as impact loss.
  This planar convention is shared with the reference implementation
  and with the Stronge model below.
\item \textbf{Contact geometry.} $\sin\theta_a = h/R - 1$ with the WPA
  nose height $h = 1.27R$, per Eq.~\eqref{eq:cushion-geom} --- not the
  snooker $7R/5$ of \citet{mathavan2010cushion}.
\end{enumerate}

\subsubsection{Stronge tangential-compliance model (stronge\_realistic)}
\label{sec:appendix-eq-stronge}

As implemented at kernel \code{0.1.0+fb3425ad}, the compliant
alternative is the planar rough-impact model of
\citet[ch.~5]{stronge2018impact}: the contact point connects through
independent normal and tangential springs, both phases are simple
harmonic, and the impact is closed-form except for locating stick--slip
transitions. The formulation reduces to a collinear problem in the
initial tangential and normal contact-point velocities
$(v_{t0} \le 0,\ v_{n0} < 0)$, with $v_{t0}$ measured against a unit
tangent chosen opposite the slip.

\paragraph{Nondimensionalization} Only the frequency ratio is
physical. Time is scaled so the normal contact frequency is
$\omega_n = 1$; compression ends at $t_c = \pi/2$ and the impact at
$t_f = (1 + e_n)\,t_c$, where $e_n = e(-v_{n0})$ from
Eq.~\eqref{eq:ec-kinematic}, additionally clamped to $[0.01, 1]$ (a
non-physical configuration must degrade, not overflow). The
restitution-phase angle is $\phi(t) = t/e_n + \tfrac{\pi}{2}(1 -
1/e_n)$, and the normal channel closes with
$v_{nf} = -e_n v_{n0}$. The tangential spring oscillates at
$\omega_t = \omega_r$ (the \emph{spin ratio} of
Table~\ref{tab:models}: the tangential-to-normal contact-stiffness
frequency ratio, equivalent to a cushion Poisson-ratio
parameterization; in-domain $\omega_r \in (1,2)$, frozen at
$1.8$), with the inverse-inertia coefficients of a sphere on a
half-space $\beta_t = 3.5$, $\beta_n = 1$,
$\beta = \beta_t/\beta_n$, and $\eta^{2} = \beta/\omega_r^{2}$.
Friction is $\mu = \mu_w(|v_{t0}|)$ from Eq.~\eqref{eq:muw-hook}.

\paragraph{Branch structure} With $r = v_{t0}/v_{n0} \ge 0$:
\begin{equation}
\begin{aligned}
&\text{initial stick:} && r < \mu\,\eta^{2},\\
&\text{gross slip:} && r > \max\bigl(\mu\beta,\;
\mu[(1+e_n)\beta - \eta^{2}/e_n]\bigr),\\
&\text{slip--stick--slip:} && \text{otherwise.}
\end{aligned}
\label{eq:stronge-branches}
\end{equation}
The gross-slip branch closes in one line,
$v_{tf} = v_{t0} - \mu\,v_{n0}\,\beta\,(1+e_n)$ (the Coulomb
limit). For initial slip, the stick onset $t_s$ has the closed form
$t_s = t_c\,T$ with
\begin{equation}
T =
\begin{cases}
\dfrac{2}{\pi}\arccos\!\dfrac{(r/\mu) - \beta}{\eta^{2} - \beta}, &
r \le \mu\beta,\\[8pt]
\dfrac{2 e_n}{\pi}\left[\arccos\!\dfrac{(r/\mu) - \beta}
{\eta^{2}/e_n - e_n\beta} - \dfrac{\pi}{2}\Bigl(1 -
\dfrac{1}{e_n}\Bigr)\right], &
r > \mu\beta .
\end{cases}
\label{eq:stronge-tstick}
\end{equation}
(onset in compression and in restitution, respectively).

\paragraph{Stick evolution and the cone-crossing functions} During
stick the tangential spring holds the contact ($v_t$ and the spring
displacement evolve harmonically at $\omega_t$). For the
\emph{initial-stick} branch the spring loads from rest at $t = 0$,
with displacement $u(t) = -(v_{t0}/\omega_t)\sin(\omega_t t)$; for the
\emph{slip--stick--slip} branch the harmonic stick episode starts at
$t_s$ from the state accumulated over the initial slip,
\begin{equation}
(u_0,\, v_0) =
\begin{cases}
\bigl(-\mu\,v_{n0}\,\eta^{2}\sin t_s,\;\;
v_{t0} - \mu\beta\,v_{n0}\,(1 - \cos t_s)\bigr), & t_s \le t_c,\\[4pt]
\bigl(-\mu\,v_{n0}\,\eta^{2}\sin\phi(t_s),\;\;
v_{t0} - \mu\beta\,v_{n0}\,(1 - e_n\cos\phi(t_s))\bigr), & t_s > t_c,
\end{cases}
\label{eq:stronge-stickstate}
\end{equation}
with $u_0$ the spring displacement and $v_0$ the tangential contact
velocity at onset. Slip resumes when the required spring force first
crosses the Coulomb cone $\mu f_n(t)$, whose restitution-phase profile
is proportional to $\sin\phi(t)$; the root-searched cone-violation
functions are, for the initial-stick branch,
\begin{equation}
f_{\mathrm{IS}}(t) \;=\; \bigl|\,r\sin(\omega_t t)\,\bigr|
\;-\; \mu\,\eta^{2}\,\omega_t\,\sin\phi(t),
\label{eq:stronge-cone-is}
\end{equation}
and for the slip--stick--slip branch, in terms of the normalized
spring state $\alpha(t)$,
\begin{equation}
\alpha(t) = \frac{u_0\cos\bigl(\omega_t(t - t_s)\bigr)
- (v_0/\omega_t)\sin\bigl(\omega_t(t - t_s)\bigr)}{\mu\,v_{n0}},
\label{eq:stronge-alpha}
\end{equation}
\begin{equation}
f_{\mathrm{SSS}}(t) \;=\; |\alpha(t)| \;-\; \eta^{2}\sin\phi(t),
\label{eq:stronge-cone-sss}
\end{equation}
a cone violation being $f > 0$. At the terminal slip time
$t_{\mathrm{slip}}$ (the first-crossing selection below) the exit
velocity combines the harmonic stick state with the Coulomb friction
accumulated over the remaining normal impulse: with
$d = t_{\mathrm{slip}} - t_s$,
\begin{align}
v_{tf}^{\mathrm{IS}} &= v_{t0}\cos(\omega_t\,t_{\mathrm{slip}})
+ \sigma\,\mu\,v_{n0}\,\beta\,e_n\,
\bigl[1 + \cos\phi(t_{\mathrm{slip}})\bigr],
\label{eq:stronge-vtf}\\
v_{tf}^{\mathrm{SSS}} &= \omega_t\,u_0\sin(\omega_t d)
+ v_0\cos(\omega_t d)
+ \sigma\,\mu\,v_{n0}\,\beta\,e_n\,
\bigl[1 + \cos\phi(t_{\mathrm{slip}})\bigr],
\label{eq:stronge-vtf-sss}
\end{align}
(the initial-stick and slip--stick--slip exits, respectively),
where the sign $\sigma = \pm 1$ is set by the side of the spring at
the crossing (friction on resumed slip opposes the slip the spring
force drives): $\sigma = +1$ iff
$-r\sin(\omega_t\,t_{\mathrm{slip}}) \ge 0$ for the initial-stick
branch and iff $\alpha(t_{\mathrm{slip}}) \ge 0$ for the
slip--stick--slip branch (in-domain, $e_n\omega_t > 1$, every first
crossing lands on $\sigma = +1$). If no crossing exists, stick holds
to $t_f$, where $\phi(t_f) = \pi$ and the friction term vanishes
identically, leaving the pure harmonic value.

\paragraph{First-crossing selection rule} At stick onset the
cone-crossing function has a structural double root
($f(t_s) = f'(t_s) = 0$ exactly: the spring rides the Coulomb limit at
onset), so an unrestricted bracketing search can lock onto the
tangency instead of the genuine crossing---flipping between root
families at float-noise scale along an incidence sweep. This was
the branch-selection defect (the $0.85$-diamond bank-arrival
discontinuity; case study
in Section~\ref{sec:oracle}). The shipped selection rule
(Section~\ref{sec:simulator-robustness}) scans the branch's
cone-violation function---$f_{\mathrm{IS}}$ of
Eq.~\eqref{eq:stronge-cone-is} or $f_{\mathrm{SSS}}$ of
Eq.~\eqref{eq:stronge-cone-sss}---over $K = 64$ equally spaced samples
on $(t_{\mathrm{lo}}, t_f]$, with $t_{\mathrm{lo}} = t_c$ for the
initial-stick branch and $t_{\mathrm{lo}} = \max(t_s, t_c)$ for the
slip--stick--slip branch; the first sample with a positive cone
violation closes a bracket containing a simple root, polished by a
deterministic Newton-plus-bisection iteration with a fixed
64-iteration cap; if no sample is positive, stick holds to $t_f$. $K$
is fixed, never adaptive. The pre-fix whole-bracket selection is
retained only as an explicit legacy option for byte-exact replay.

\paragraph{Impulse application} The contact-velocity changes map to
the ball through the inverse-inertia coefficients,
\begin{equation}
\Delta v_n = \frac{v_{nf} - v_{n0}}{\beta_n}, \qquad
\Delta v_t = \frac{v_{tf} - v_{t0}}{\beta_t}, \qquad
\Delta\boldsymbol{\omega} = \frac{5}{2R}\,
(-\hat{\mathbf{n}})\times\hat{\mathbf{t}}\;\Delta v_t,
\label{eq:stronge-apply}
\end{equation}
where $\hat{\mathbf{n}}$ is the three-dimensional contact normal
tilted below the horizontal by $\theta_a$ from
Eq.~\eqref{eq:cushion-geom} (a deviation from the reference
implementation's $\tan\theta_a$ convention, worth about $0.5^\circ$),
and the outgoing $v_z$ is zeroed (planar convention, as in the Han
model).

\subsubsection{Mathavan dual-contact model and the fitted presets
(mathavan\_fitted, grip\_fitted)}
\label{sec:appendix-eq-mathavan}

As implemented at kernel \code{0.1.0+fb3425ad}, the fitted presets are
built on the dual-contact impulse-space integration of
\citet{mathavan2010cushion}: the impulse accumulates in steps
$\Delta P$ while Coulomb friction acts at \emph{two} contact points,
the cushion nose I (friction $\mu_w$) and the table bed C (friction
$\mu_s$, the cloth sliding coefficient), the bed normal force growing
as the nose pushes the ball down. In the cushion frame
($\hat{\mathbf{x}}$ along the rail, $\hat{\mathbf{y}}$ the horizontal
normal into the cushion, approach $v_y > 0$), the contact-point slip
velocities are
\begin{equation}
\begin{aligned}
\text{nose I:}\quad &v_{xI} = v_x + \omega_y R\sin\theta_a -
\omega_z R\cos\theta_a, &\quad
&v_{yI} = -v_y\sin\theta_a + \omega_x R,\\
\text{bed C:}\quad &v_{xC} = v_x - \omega_y R, &
&v_{yC} = v_y + \omega_x R,
\end{aligned}
\label{eq:mathavan-slips}
\end{equation}
with slip angles $\phi_I, \phi_C$ and nose force ratios
$(g_x, g_y) = \mu_w(\cos\phi_I, \sin\phi_I)$ on the pure-Coulomb path.
One impulse step updates (the paper's eqs.~12--14, transcribed as
implemented; $D = \max(\sin\theta_a + g_y\cos\theta_a,\, 0)$ is the
bed normal-force factor)
\begin{equation}
\begin{aligned}
\Delta v_x &= -\tfrac{1}{m}\bigl(g_x +
\mu_s\cos\phi_C\,D\bigr)\Delta P,\\
\Delta v_y &= -\tfrac{1}{m}\bigl(\cos\theta_a - g_y \sin\theta_a +
\mu_s\sin\phi_C\,D\bigr)\Delta P,\\
\Delta \omega_x &= -\tfrac{5}{2mR}\bigl(g_y +
\mu_s\sin\phi_C\,D\bigr)\Delta P,\\
\Delta \omega_y &= -\tfrac{5}{2mR}\bigl(g_x\sin\theta_a -
\mu_s\cos\phi_C\,D\bigr)\Delta P,\\
\Delta \omega_z &= +\tfrac{5}{2mR}\,g_x\cos\theta_a\,\Delta P .
\end{aligned}
\label{eq:mathavan-step}
\end{equation}
Compression integrates until $v_y \le 0$ (the crossing re-approached
with eight fixed step-halvings), accumulating the nose normal work
$W_{\mathrm{comp}}$ via $\mathrm{d}W = \Delta P\,|v_y|\cos\theta_a$;
restitution integrates until the nose work reaches the energetic
target
\begin{equation}
W_{\mathrm{rest}} \;=\; ee^{2}\,W_{\mathrm{comp}},
\label{eq:mathavan-ee}
\end{equation}
with $ee = ee(v_{y0})$ from Eq.~\eqref{eq:ec-energetic} and the final
step sized exactly to the remaining work. The step is
$\Delta P = \max(m\,v_{y0}/N,\ 10^{-4}\,\text{kg\,m/s})$ with
$N = 1000$ (restitution analogously from the target work), a hard cap
of $10N$ steps per phase raising a diagnostic error (never silent
truncation), and impacts whose normal approach exceeds the
$2.5$\,m/s experimental validity bound of the source are counted per
shot rather than rejected. Itemized deviations from the published
paper: (i) restitution enters through the energetic hook of
Eq.~\eqref{eq:ec-energetic} rather than the paper's fixed snooker fit;
(ii) nose friction enters through Eq.~\eqref{eq:muw-hook} at the
arrival nose-slip speed, bed friction is the cloth $\mu_s$; (iii) the
bed normal-force factor is clamped at zero ($D \ge 0$: the unilateral
bed cannot pull); (iv) the paper's \S3.2 rolling-condition force
nulling (its eqs.~15a/15b) is omitted---measured effect below $+0.06$
on the rolling-perpendicular exit-speed ratio---because an
epsilon-free rolling detector does not exist in impulse stepping;
(v) the incoming $v_z$ passes through unchanged (planar convention).

\paragraph{Tangential compliance (both fitted presets)} The fitted
presets add a first-party tangential-compliance element at the nose: a
Cundall--Strack-style bristle spring per nose tangent direction, fed
by the contact-point slip, its per-step force Coulomb-capped at
$\mu_w N$. The impulse-to-time map is a half-sine normal-force
profile, floored and argument-capped, over the a-priori total-impulse
estimate:
\begin{equation}
N(p) = \max\!\left[
N_{\mathrm{peak}}\sin\!\Bigl(\frac{\pi\,\bar{p}}{P_{\mathrm{tot}}}\Bigr),\;
N_{\mathrm{peak}}\sin(0.02\,\pi)\right],
\quad
\bar{p} = \min(p,\, 0.999\,P_{\mathrm{tot}}),
\label{eq:bristle}
\end{equation}
\begin{equation}
N_{\mathrm{peak}} = \frac{\pi\,P_{\mathrm{tot}}}{2\,\tau_c},
\quad
P_{\mathrm{tot}} = \frac{m\,v_{y0}\,(1 + ee)}{\cos\theta_a},
\quad
k_t = \Bigl(\frac{\chi\pi}{\tau_c}\Bigr)^{2}\frac{m}{\beta_t},
\label{eq:bristle-kt}
\end{equation}
where $p$ is the accumulated normal impulse and the integrator
advances internal time as $\mathrm{d}t = \Delta P / N(p)$. The
$\sin(0.02\pi)$ floor and the $0.999\,P_{\mathrm{tot}}$ argument cap
are part of the frozen arithmetic, not implementation detail: the bare
profile gives $N(0) = 0$, making the first step's $\mathrm{d}t$
singular, and both guards reshape the profile tails. The internal time
scale $\tau_c$ cancels identically (only the frequency ratio $\chi$ is
physical, the same invariance as the Stronge normalization) and
$\beta_t = 3.5$. The peak elastic energy stored in
the follow-channel spring is held as a reservoir and returned at
release as an impulse along the stored displacement with gain $e_x$
(energy returned $= e_x^{2} \times$ held), capped by both the
whole-impact kinetic-energy deficit and the preset-dependent Coulomb
allowance $\mu_{\mathrm{rel}}\,(P_{\mathrm{tot}} -
P_{\mathrm{comp}})$, with $P_{\mathrm{comp}}$ the normal impulse
accumulated at the end of compression. The release ceiling is the
\emph{static} grip bound for a captured \code{grip_fitted}
impact---$\mu_{\mathrm{rel}} = \mu_{\mathrm{grip}} = 2.5\,\mu_w$ (the
held reservoir is grip-phase elastic storage; its give-back flows
through the gripped interface, not a sliding one)---and the kinetic
$\mu_{\mathrm{rel}} = \mu_w$ otherwise (\code{mathavan_fitted} and
uncaptured arrivals); $e_x$ is a
phenomenological give-back gain, not a Cross-convention tangential
restitution. \code{mathavan_fitted} freezes $\chi = 0.8$,
$e_x = 0.9$ (jointly fitted with $ee_0 = 0.80$ and the $\mu_w$ clamps
above against a high-speed-video response surface
\citep{alciatore_hsv_catalog}, with the source's own rebound figures
held as anchors). \code{grip_fitted} additionally enables a
static/kinetic stick--slip mode per nose channel: the loaded spring
transmits force up to the static ceiling
$\mu_{\mathrm{grip}} = 2.5\,\mu_w$ while gripping, engaged only for
spin-overdriven arrivals with nose slip at or below
$0.6$\,m/s (capture threshold), with channel disengage when the spring
unloads through its natural length and a whole-impact energy ledger
that scales tangential exit degrees of freedom down if the loop would
otherwise exit above the arrival kinetic energy. Its constants are
$\chi = 1.2$, $e_x = 0.9$, grip ratio $2.5$, capture $0.6$\,m/s,
paired with Eq.~\eqref{eq:ec-grip}. Both presets' declared validity
domains and out-of-domain observables are ledgered;
neither was re-tuned.

\subsection{Compliant frozen-cluster resolution (CompliantChain)}
\label{sec:appendix-eq-cluster}

As implemented at kernel \code{0.1.0+fb3425ad}, an impact reaching a
frozen cluster (any chain of contacts within $10^{-5}$\,m surface
separation, swept transitively) escalates to a bounded compliant
episode (Section~\ref{sec:simulator-break}). Every active contact
carries the Hertz--Lankarani--Nikravesh normal force
\citep{lankarani1990contact,lankarani1994continuous}
\begin{equation}
F_N \;=\; \delta^{3/2}\,\bigl(K + \chi_{\mathrm{LN}}\,\dot\delta\bigr)
\;\ge\; 0,
\qquad
\chi_{\mathrm{LN}} \;=\; \frac{3K\,(1 - e_b^{2})}
{4\,\max(\dot\delta^{-},\,10^{-3}\,\text{m/s})},
\label{eq:ln-force}
\end{equation}
with $\delta$ the penetration depth, $e_b$ the mean ball--ball
restitution of the pair, and $\dot\delta^{-}$ the contact's impact
speed, maintained as the largest \emph{equivalent} impact speed
$\smash{\sqrt{\dot\delta^{2} + 4K\delta^{5/2}/(5 m_{\mathrm{eff}})}}$
the contact has exhibited (the Hertz energy inversion; for clean
impacts this equals the first-penetration approach speed exactly, and
for contacts entering a window mid-compression it restores the
Lankarani--Nikravesh definition of $\dot\delta^{-}$ as the impact
velocity). The force is clamped non-negative (cohesionless contact),
and the grazing clamp bounds the damping singularity as
$\dot\delta^{-} \to 0$. The Hertz stiffness is derived, not fitted:
\begin{equation}
K = \tfrac{4}{3}\,E^{*}\sqrt{R^{*}}, \qquad
\frac{1}{E^{*}} = \frac{2\,(1-\nu^{2})}{E}, \qquad
R^{*} = \frac{R}{2},
\label{eq:hertz-k}
\end{equation}
with $E = 3800$\,MPa and $\nu = 0.375$ for phenolic resin
\citep[Table~1]{gismerosmoreno2022multiple} and the catalog ball
radius, giving $K = 3.524\times10^{8}$\,N/m$^{3/2}$. Tangential
traction is force-based Coulomb, $\mu_{bb}(v_{\mathrm{slip}})\,F_N$
opposing the contact-point slip with the same friction curve as
Eq.~\eqref{eq:bb-mu} (no impulse cap is consumed---the property that
kept the compliant rows bit-stable across the mid-program re-baseline,
Section~\ref{sec:simulator-break}).

All contacts advance \emph{concurrently} (Jacobi accumulation: every
force reads only previous-step state, making the episode order-free
and mirror-equivariant) under semi-implicit Euler with fixed step
$\mathrm{d}t = 10^{-6}$\,s and a hard cap of $3000$ steps (a 3\,ms
window bounding the 1--2\,ms break transient). The episode ends early
when every contact is open ($\delta \le 0$) and non-approaching
(normal closing speed below $10^{-9}$\,m/s), or hands off mid-window
when a non-member pair first closes within $10^{-5}$\,m while
approaching (the event loop detects the new contact and re-escalates
with the newcomer as a member). Three post-window guards fail closed,
restoring the pre-window state and running the legacy resolver on the
same inputs: non-finite state; cluster energy gain (kinetic plus
stored Hertz elastic energy $\tfrac{2}{5}K\delta^{5/2}$ exceeding the
pre-window value by more than $10^{-9}$\,J); and cap expiry with a
still-penetrating approaching pair. The \code{SequentialImpulse}
legacy arm resolves the same cluster by deterministic Gauss--Seidel
sweeps of pairwise frictional impulses (at most 256 sweeps) and is
retained as an explicit recorded option; it is the baseline class
whose starved momentum exchange produces the dead-cue break defect
that experiment \expid{E3} gates on
(Section~\ref{sec:simulator-break}).

\subsection{Event detection, epsilons, and rest projection}
\label{sec:appendix-eq-events}

As implemented at kernel \code{0.1.0+fb3425ad}, every candidate event
time is a real root of a polynomial in time of degree at most four,
because every ball-position component is at most quadratic
(\ref{sec:appendix-eq-cloth})---with one \emph{scheduled} exception,
the airborne pocket heuristic below. Per event class: ball--ball
proximity $|\Delta\mathbf{r}(\tau)|^{2} = (2R)^{2}$ between two
quadratic paths is quartic; a linear cushion (vertical plane at signed
distance $R$) is quadratic; jaw arcs and pockets (vertical cylinders,
contact at arc radius $+R$; capture of a \emph{grounded} ball when the
\emph{center} crosses the pocket radius) are quartic; the regime
transitions have the closed
forms $\tau_{\mathrm{slide}}$, $\tau_{\mathrm{roll}}$,
$\tau_{\mathrm{spin}}$ of \ref{sec:appendix-eq-cloth} and
need no root-finding. Quartics are solved with the closed-form solver
of \citet{orellana2020algorithm1010} incorporating both published
corrections \citep{demichele2022remark,demichele2024remark}, plus two
documented numerical accommodations inherited from
\citet{kiefl2024pooltool}: the near-degenerate-quadratic threshold in
the LDLT factorization is loosened by a factor of 100 so that
Newton's-cradle-like quartics (nearly equal velocities) yield four
distinct roots rather than spurious doubles, and a root is accepted as
real under a hybrid filter---$|\mathrm{Im}| < 10^{-9}$ when
$|\mathrm{Re}| > 10^{-3}$, else $|\mathrm{Im}|/|\mathrm{Re}| <
10^{-3}$. The earliest accepted root strictly beyond the $10^{-9}$\,s
simultaneity guard becomes the next event; a leading coefficient is
degraded (quartic to cubic, and so on) when it cannot exceed one
machine epsilon of the next term anywhere within the Cauchy root
bound. Contact at $\tau = 0$ fires immediately only for approaching
pairs; an at-contact cushion root requires approach speed above
$10^{-5}$\,m/s (the sleeping threshold that lets resting contact
coexist with detection).

\paragraph{Airborne pocket capture (non-root scheduling)} For an
\emph{airborne} ball, pocket capture is not a polynomial root: it is a
documented phenomenological approximation inherited from
\citet{kiefl2024pooltool}. With $\tau_{\mathrm{land}}$ the ballistic
descending crossing of $z = R$: (i) if the landing point lies strictly
inside the pocket radius, capture is scheduled at
$\tau_{\mathrm{land}} - \epsilon_t$ (the $10^{-9}$\,s time epsilon),
guaranteeing temporal precedence over the bed-landing event; (ii)
otherwise the horizontal flight path (linear in $xy$) is intersected
with the pocket cylinder, giving quadratic entry/exit roots
$\tau_{\mathrm{in}} < \tau_{\mathrm{out}}$, and capture is scheduled
at the midpoint $(\tau_{\mathrm{in}} + \tau_{\mathrm{out}})/2$ iff
$\tau_{\mathrm{in}} \ge 0$, $z(\tau_{\mathrm{in}}) \ge R$ (the ball
did not land first), and $z(\tau_{\mathrm{out}}) \le \tfrac{7}{5}R$
(the rim heuristic: low fliers ricochet off the back of the pocket
and drop in, higher ones fly over); otherwise no airborne capture
candidate exists and the landing proceeds. A jump-shot reproduction
that treats pocket capture as a root of the quartic above will
mis-order these events.

Frozen tolerance table: time, space, velocity, and spin epsilons all
$10^{-9}$ (s, m, m/s, rad/s); root-acceptance bounds as above;
separation (``kiss'') gap $10^{-6}$\,m; frozen-rack surface gap
$0.5\,\mu$m (below the $10^{-5}$\,m cluster threshold, so racked
contacts are frozen topology; seeded back-row micro-gaps
0.05--0.5\,mm sit above it).

\paragraph{Frozen-rack topology (seeded selection rule)} The default
rack constructor realizes the frozen topology with frozen fraction
$f = 0.7$ (a declared first-party consensus estimate; no published
frozen-fraction statistics exist). The apex wedge---apex plus row
1---always keeps its lattice nodes and consumes no random draw. Each
back-row ball (lattice row $\ge 2$), visited in rack order against the
shot-seeded deterministic stream, keeps its node with probability $f$;
otherwise it is displaced by a micro-gap $g \sim U[0.05, 0.5)$\,mm
along a direction $\vartheta \sim U[0, 2\pi)$, clamped so no pair
drops below the frozen lattice pitch. Then each back-row boundary
$r{-}1 \to r$ ($r \ge 2$, ascending, same stream) opens with
probability $1 - f$ by shifting all rows $\ge r$ away from the breaker
by an independent $g$. The same seed therefore regenerates the exact
rack topology of experiment \expid{E3} bit-for-bit; the fraction and
rule are required inputs, not tuning detail.

\paragraph{Terminal rest projection (round-trip law)} Every
simulation ends by projecting all balls to exact rest, and any pair
the input validator would reject---center distance below
$r_a + r_b - 10^{-6}$\,m---is retracted along the line of centers,
each ball moving half the correction, to the exact $10^{-6}$\,m gap,
sweeping pairs in ball-identity order to a fixed point (at most 128
passes; observed convergence in single-digit passes), then re-seated
to $z = R$ and clamped against rail-seat bands and pocket capture
discs. Pairs the validator accepts, including undisturbed frozen-rack
contacts, are never touched, so clean records remain byte-identical.
The result is the round-trip law of
Section~\ref{sec:simulator-robustness}: every \code{simulate} output
is a valid \code{simulate} input, with no accumulating
representational debt over arbitrarily long shot chains.

\subsection{Consolidated frozen constants}
\label{sec:appendix-eq-constants}

Table~\ref{tab:frozen-constants} consolidates every constant fixed by
the frozen catalog, each with units and provenance class:
\emph{lit} (published value or derived from published values),
\emph{fit} (first-party calibration, declared as such), or
\emph{num} (numerical guard or discretization constant). All values
are as implemented at kernel \code{0.1.0+fb3425ad} and were frozen
before any battery sweep ran (Section~\ref{sec:methodology}).

{\footnotesize
\begin{longtable}{@{}>{\raggedright\arraybackslash}p{0.31\linewidth}>{\raggedright\arraybackslash}p{0.22\linewidth}c>{\raggedright\arraybackslash}p{0.27\linewidth}@{}}
  \caption{Consolidated frozen constants of the validated models.
  Provenance classes: lit = published/derived, fit = first-party
  calibration (declared), num = numerical guard. Equation references
  are to this appendix.}
  \label{tab:frozen-constants}\\
  \toprule
  Constant & Value & Class & Provenance / role \\
  \midrule
  \endfirsthead
  \toprule
  Constant & Value & Class & Provenance / role \\
  \midrule
  \endhead
  \bottomrule
  \endfoot
  \multicolumn{4}{@{}l}{\emph{Ball, cloth, and gravity
    (\ref{sec:appendix-eq-cloth})}}\\
  Ball mass $m$ & $0.170097$\,kg & lit & regulation set
    \citep{kiefl2024pooltool} \\
  Ball radius $R$ & $0.028575$\,m & lit & $57.15$\,mm US ball \\
  Gravity $g$ & $9.81$\,m/s$^2$ & lit & --- \\
  Sliding friction $\mu_s$ & $0.20$ & lit & worsted tournament cloth
    \citep{alciatore_faq} \\
  Rolling friction $\mu_r$ & $0.010$ & lit & lag-speed calibration
    class (Section~\ref{sec:results}) \\
  Spin-decay $\mu_{sp}$; rate $\alpha$ & $\tfrac{4}{9}\,(R/\mathrm{m})
    = 0.0127$ (dimensionless);
    $\tfrac{10}{9}g \approx 10.9$\,rad/s$^2$ & lit &
    video-measured $\approx 10$\,rad/s$^2$
    \citep{alciatore_tp_catalog}; Eq.~\eqref{eq:spin-decay} \\
  Ball--slate restitution $e_t$ & $0.5$ & lit & reference default
    \citep{kiefl2024pooltool}; no published measurement;
    Eq.~\eqref{eq:slate-normal} \\
  Slate bounce guard (apex) & $5$\,mm (default) / $10\,\mu$m (realism
    presets) & num & Zeno guard, \ref{sec:appendix-eq-slate} \\
  \addlinespace
  \multicolumn{4}{@{}l}{\emph{Ball--ball impact
    (\ref{sec:appendix-eq-ballball})}}\\
  Restitution $e_b$ & $0.95$ & lit & community constant:
    \citet{alciatore_faq} physical-properties compilation (band
    0.92--0.98, printed without attribution; no first-hand measurement
    located; the $[0.92, 0.98]$
    bracket is a declared sensitivity range, not a measured interval
    --- main-text Section~4); ledgered as declared model content \\
  $\mu_{bb}(v)$: $(a, b, c)$ & $(0.009951,$ $0.108,$ $1.088)$ & lit &
    \citet{alciatore_tp_a14} fit; Eq.~\eqref{eq:bb-mu} \\
  Slip-impulse cap & $(1+e_b)/2 \cdot |\Delta v_n^{-}|$ & lit &
    TP~A.5 Eq.~10; corrected default (slip-law fix) \\
  Legacy cap (replay only) & $e_b\,|\Delta v_n^{-}|$ & lit &
    pre-correction law \citep{kiefl2024pooltool} \\
  Micro-hop keep threshold & $5$\,mm (default) / $10\,\mu$m
    (realism presets) & num & bounce-apex guard,
    \ref{sec:appendix-eq-ballball} \\
  \addlinespace
  \multicolumn{4}{@{}l}{\emph{Stick--ball strike
    (\ref{sec:appendix-eq-strike})}}\\
  Cue mass $M$ & $0.567$\,kg & lit & 20\,oz cue \\
  Tip radius $\rho$ & $0.0106$\,m (nickel); $0.00895$\,m dime variant
    & lit & TP~B.1 curvature standards \\
  Default endmass $m_e$ & $m/30 \approx 5.67$\,g & lit & TP~A-31
    \citep{alciatore_tp_a31} \\
  Endmass catalog span & $1.70$--$12.7$\,g & lit & pivot-length
    inversion $m/m_e = 2.5 L_p/R - 3.5$
    \citep{shepard2001squirt}; robot-measured $L_p$ \\
  Tip restitution $e_{\mathrm{tip}}$ & $1.0$ default; $0.73$ leather;
    $0.85$ phenolic & lit & measured tip efficiency
    \citep{alciatore_hsv_catalog} \\
  Miscue envelope & $a_{\mathrm{cmd}}^2 + b_{\mathrm{cmd}}^2 \le 0.25$
    & lit & tip-offset limit, Eq.~\eqref{eq:strike-miscue};
    contact-offset binding post-registration
    (\ref{sec:appendix-eq-strike}) \\
  Side-offset sign map & $a = -a_{\mathrm{cmd}}/(1+\rho/R)$ & num &
    commanded (right-English positive) $\to$ internal (left positive);
    Eq.~\eqref{eq:strike-tipscale} \\
  Shot speed domain & $v_0 \in (0, 15]$\,m/s & num & input validity \\
  Pinch onset / full & $60^{\circ}$ / $70^{\circ}$ & lit &
    jump-vs-mass\'e regime boundary \citep{coriolis1835};
    Eq.~\eqref{eq:strike-pinch} \\
  Press return $\kappa$; relax $\tau_{\mathrm{pinch}}$ & $0.45$;
    $0.12$\,s & fit & mass\'e footage fit;
    Eq.~\eqref{eq:strike-press} \\
  \addlinespace
  \multicolumn{4}{@{}l}{\emph{Cushion geometry and hooks
    (\ref{sec:appendix-eq-cushion})}}\\
  Nose height $h$ (US pool) & $1.27R$ ($\sin\theta_a = 0.27$) & lit &
    WPA specification; Eq.~\eqref{eq:cushion-geom} \\
  Nose height $h$ (snooker pairing) & $1.4R$ ($\sin\theta_a = 0.4$) &
    lit & \citet{mathavan2010cushion} \\
  Kinematic $e_0$; clamp & $0.85$; $[0.6, 0.98]$ & lit &
    \citet{han2005dynamics} lineage; Eq.~\eqref{eq:ec-kinematic} \\
  Energetic $ee_0$ (default / \code{mathavan_fitted}) & $0.98$ /
    $0.80$ & lit / fit & snooker fit \citep{mathavan2010cushion} /
    joint video fit \\
  \code{grip_fitted} $ee(v)$ & $0.65 + 0.18(v_n - 1)$, clamp
    $[0.5, 0.98]$ & fit & Eq.~\eqref{eq:ec-grip} \\
  $\mu_w$ (\code{han2005_default}) & $0.20$ flat & lit & no published
    $\mu_w(v)$ exists (declared) \\
  $\mu_w$ (\code{stronge_realistic}) & $0.20 - 0.4(v-1)$, clamp
    $[0.20, 0.40]$ & fit & video kick/sidespin anchors;
    Eq.~\eqref{eq:muw-hook} \\
  $\mu_w$ (fitted presets) & $0.20 - 0.4(v-1)$, clamp $[0.16, 0.28]$
    & fit & joint fit with tangential compliance \\
  \addlinespace
  \multicolumn{4}{@{}l}{\emph{Stronge model
    (\ref{sec:appendix-eq-stronge})}}\\
  Spin ratio $\omega_r$ & $1.8$; domain $(1, 2)$ & lit &
    \citet{stronge2018impact} parameterization; reference default
    \citep{kiefl2024pooltool} \\
  Inverse inertia $\beta_t$, $\beta_n$ & $3.5$, $1$ & lit & sphere on
    half-space \citep{stronge2018impact} \\
  $e_n$ clamp & $[0.01, 1]$ & num & degradation guard \\
  First-crossing scan $K$; root cap & $64$; $64$ iterations & num &
    fixed-resolution selection (branch-selection fix) \\
  \addlinespace
  \multicolumn{4}{@{}l}{\emph{Mathavan model and fitted presets
    (\ref{sec:appendix-eq-mathavan})}}\\
  Impulse step $\Delta P$ & $\max(m v_{y0}/1000,$ $10^{-4})$ kg\,m/s
    & num & first-order bias control \\
  Step caps; refinement & $10\times1000$ per phase; $8$ halvings &
    num & diagnostic overrun, never silent \\
  Validity bound & $v_n \le 2.5$\,m/s (counted) & lit & source
    experimental range \citep{mathavan2010cushion} \\
  Bristle force floor; cap & $N \ge N_{\mathrm{peak}}\sin(0.02\pi)$;
    $p \le 0.999\,P_{\mathrm{tot}}$ & num & profile-tail guards;
    Eq.~\eqref{eq:bristle} \\
  \code{mathavan_fitted} $\chi$, $e_x$ & $0.8$, $0.9$ & fit &
    joint HSV-response fit; Eq.~\eqref{eq:bristle-kt} \\
  \code{grip_fitted} $\chi$, $e_x$, grip ratio, capture & $1.2$,
    $0.9$, $2.5$, $0.6$\,m/s & fit & grip-mode joint fit \\
  \addlinespace
  \multicolumn{4}{@{}l}{\emph{Break cluster
    (\ref{sec:appendix-eq-cluster})}}\\
  Hertz stiffness $K$ & $3.524\times10^{8}$\,N/m$^{3/2}$ & lit &
    derived, Eq.~\eqref{eq:hertz-k}
    \citep{gismerosmoreno2022multiple} \\
  Step; cap & $10^{-6}$\,s; $3000$ steps (3\,ms) & num & Hertz
    contact-time resolution \\
  Grazing clamp on $\dot\delta^{-}$ & $10^{-3}$\,m/s & num &
    Lankarani--Nikravesh singularity guard; Eq.~\eqref{eq:ln-force} \\
  Cluster contact threshold & $10^{-5}$\,m & num & escalation /
    hand-off tolerance \\
  Energy-audit slack & $10^{-9}$\,J & num & fail-closed guard \\
  Legacy sweep cap & $256$ & num & \code{SequentialImpulse} bound \\
  \addlinespace
  \multicolumn{4}{@{}l}{\emph{Event detection and rest
    (\ref{sec:appendix-eq-events})}}\\
  Time/space/velocity/spin eps & $10^{-9}$ (s, m, m/s, rad/s) & num &
    simultaneity and zero guards \\
  Root-real filter & $|\mathrm{Im}| < 10^{-9}$ if $|\mathrm{Re}| >
    10^{-3}$, else ratio $< 10^{-3}$ & num & hybrid acceptance \\
  Degeneracy relaxation & $\times 100$ & num & cradle-quartic guard
    \citep{kiefl2024pooltool} \\
  Separation gap & $10^{-6}$\,m & num & retraction target;
    round-trip law \\
  Frozen-rack gap; seeded gaps & $0.5\,\mu$m; $0.05$--$0.5$\,mm &
    num & frozen-contact topology
    (Section~\ref{sec:simulator-break}) \\
  Frozen fraction $f$ (rack) & $0.7$; apex wedge always frozen,
    back rows sampled & fit & racking-consensus estimate (declared);
    selection rule in \ref{sec:appendix-eq-events} \\
  Rest sleeping threshold & $10^{-5}$\,m/s & num & at-contact
    approach gate \\
  Rest projection pass cap & $128$ & num & fixed-point sweep bound \\
\end{longtable}
}

% ===========================================================================
% Supplement — extended per-experiment results for groups B and D (moved
% from the main text's Sections 5.2 and 5.4 when the main text was
% shortened; the main text keeps the group summary tables and the
% non-obvious findings). Every number traces to the study's baseline
% manifest, its ledger, or a field of a pinned group-B or group-D output
% (canonical kernel 0.1.0+fb3425ad). Figure and table references without
% an "S" prefix resolve to the main document (xr).
% ===========================================================================
\section{Extended results: groups B and D}
\label{sec:supp-extended-results}

This section carries the per-experiment result prose behind the group
summary tables of the main text's Sections~\ref{sec:results-b}
(stick--ball) and~\ref{sec:results-d} (curved paths). The figures cited
here (Figs.~\ref{fig:F6} and~\ref{fig:F10}) and the miscue-convention
table (Table~\ref{tab:miscue-conventions}) appear in the main document.

\subsection{Group B --- stick--ball, per experiment}
\label{sec:supp-results-b}

\paragraph{Endmass squirt (\expid{B1})}
The kernel's cue-ball deflection (squirt) is the endmass closed form of
TP A-31 \citep{alciatore_tp_a31} and \citet{cross2008squirt}, in the
lineage of \citet{shepard2001squirt}. Over contact offsets $0$--$0.5R$,
end masses $\{5, 10, 15\}$\,g, and speeds $\{1, 2, 4\}$\,m/s plus the
six-cue catalog grid (333 runs), the maximum comparison error against
the closed form is $8 \times 10^{-15}$ degrees --- machine precision,
confirming the implementation and the tip-offset plumbing end to end
(Fig.~\ref{fig:F6}a). The source's printed example evaluations reproduce
to $|E| \le 4.3 \times 10^{-4}$ degrees, the PDF's own rounding scale.
The registered qualitative features quantify as: near-linearity in
offset with $R^2 = 0.9967$--$0.9979$ per registered end-mass/speed
series ($R^2 \geq 0.9961$ for the plotted six-cue catalog grids of
Fig.~\ref{fig:F6}a), the local slope
density falling $\sim$12\% toward the miscue limit (mildly sub-linear);
speed invariance exact (measured spread $0.0^\circ$ --- the closed form
contains no speed); and left/right symmetry to $1.4 \times 10^{-14}$
degrees.

\paragraph{Natural pivot length (\expid{B2})}
A simulated aim-and-pivot binary search per cue recovers the natural
pivot length $L^{*}$ at which backhand-english pivoting self-corrects
squirt; geometric consistency with the closed form holds to
$4.8 \times 10^{-10}$\,m, the bisection floor. TP B.26's claim that one
pivot length serves all offsets is strongly confirmed: re-using each
cue's single $L^{*}$ across the legal offset ladder leaves a worst-case
residual aiming error of $0.077^\circ$ (shipped catalog cues
$0.006$--$0.026^\circ$). The comparison's boundary is a convention
(Fig.~\ref{fig:F6}b): catalog pivot lengths span $10.5$--$64.9$\,in,
with only the break cue inside the Platinum-protocol $8$--$14$\,in band
(the next-lowest shaft misses the band top by 0.03\,in), and the
companion squirt-angle comparison of \expid{B1} shows why --- commanding
the protocol's published 6/12\,mm as \emph{tip-center} offsets (the
protocol page's literal wording) undershoots all four robot-measured
shafts by $0.32$--$0.42^\circ$, while reading the same numbers as
\emph{contact} offsets closes the gap almost exactly ($+0.02^\circ$ on
the lowest-squirt shaft) with a growing positive residual toward the
break cue ($+0.22^\circ$). Re-referencing the pivot geometry to contact
offsets scales $L^{*}$ by $1/(1 + r_{\mathrm{tip}}/R) \approx 0.729$,
which would bring two further shafts inside the band. Both readings are
recorded and neither is adopted as a recalibration: the
commanded/tip-center form is the physical backhand-english geometry. In $u_{\mathrm{val}}$ terms
(Table~\ref{tab:uval}) the robot source publishes a band with no
per-shaft uncertainty, and the $0.32$--$0.42^\circ$ convention gap is an
order larger than any plausible protocol $u_D$, so the boundary here is
convention, not measurement uncertainty.

\paragraph{Miscue limit and maximum sidespin (\expid{B3})}
The registered ``near-linear'' spin-versus-offset question resolves to
exactly linear: the source relation is linear in contact offset ---
spin-rate factor $\mathrm{SRF} = \omega R / v = (5/2)(x/R)$ --- and the
kernel reproduces it to $|E| \le 2.2 \times 10^{-16}$ in SRF with fitted
quadratic curvature $\le 5.5 \times 10^{-15}$, exact odd symmetry, and
exact speed invariance (252 runs, plus a six-run dime-tip arm). The
envelope, by contrast, is a
conventions result (Table~\ref{tab:miscue-conventions};
Fig.~\ref{fig:F6}c). The kernel's miscue rule is input validation on
\emph{commanded tip-center} offsets, $a^2 + b^2 \le (0.5R)^2$, with no
tip-friction model anywhere; the source (TP 2.1) derives the limit from
tip friction at the \emph{contact} offset. Under the shipped nickel tip
the reachable contact offset therefore tops out at $0.365R$
($\mathrm{SRF} = 0.912$, below natural-roll-rate sidespin), and the tip
friction coefficients implied by the kernel's ceilings
($0.39$--$0.58$ across conventions) sit at or below the
$0.577$--$0.659$ the literature attributes to chalked tips --- every
accepted command is physically non-miscuing, so the conservatism is
one-sided. Contact offsets in $(0.365, 0.5]R$ are physically legal per
the source yet unreachable under the nickel tip, and $(0.5, 0.55]R$ is
unreachable under every convention. The registered $0.525$--$0.6R$
sweep tail was realized as rejection probes: all 16 cells reject with
identical diagnostics under both tip radii --- designed rejection
evidence for the envelope, not a failure mode. Post-registration, a
kernel amendment (kernel \code{0.1.0+8efd660e}) re-binds the miscue limit
to the contact offset (commanded cap
$0.5\,(1 + r_{\mathrm{tip}}/R) \approx 0.686R$, nickel); every
previously accepted input reproduces bit-exactly. The rejection probes
above are boundary probes of the registered envelope itself: under the
revised, cue-aware convention the eight nickel-tip cells (contact
$0.383$--$0.438R$, inside the $0.5R$ contact limit) and the registered
combined-axis probe at commanded $(0.375, 0.375)$ would be accepted
and simulate, the eight contact-mode cells still reject, and the
recorded tip-radius-independent diagnostics are superseded by
cue-aware validation by design. The pinned artifact stands unchanged
as the registered-envelope record; the post-registration maintenance
re-baseline under the revised envelope, which measures exactly this
disclosed class, is recorded in the study's baseline manifest.
$(0.5, 0.55]R$
remains unreachable (amendment note: main-text
Section~\ref{sec:limitations}).

\paragraph{Center of percussion (\expid{B4})}
Coriolis' center-of-percussion result --- immediate natural roll when the
tip strikes at $70\%$ of ball height, i.e., a contact offset of $0.4R$
above center \citep{coriolis1835,alciatore2008coriolis} --- verifies end
to end through the event pipeline over 41 offsets $\times$ three speeds
(123 runs; Fig.~\ref{fig:F6}d). The slide-phase duration is exactly zero
at $b = 0.4R$ at all three speeds; the interpolated crossing sits at
$b^{*} = 0.4R + 1.4 \times 10^{-11}$ (the documented contact-mode tip
correction), equivalently at height fraction 0.700, and the slide-time
V-shape's arms match the closed form. The spin-rate factor follows
$(5/2)(b/R)$ exactly (fitted slope 2.5, $R^2 = 1$). Above the crossing
the cue ball overspins and cloth friction acts \emph{forward}: the
roll-speed conversion $v_{\mathrm{roll}} = v(5 + 2\,\mathrm{SRF})/7$
holds to machine precision, giving $+7.14\%$ speed at $b = 0.5R$; below
it, a draw strike retains $35.7\%$ of strike speed at $b = -0.5R$ and
never reverses direction without an object ball. The slide map matches
the closed forms to $6.3 \times 10^{-11}$\,m in slide distance and
$1.8 \times 10^{-11}$\,s in duration over the 104 uncensored cells; the
19 censored cells are all $v_0 = 3$\,m/s near-stun or draw shots whose
closed-form slide distance (up to ${\approx}3.2$\,m) exceeds the available
$\sim$2.44\,m runway --- slides that cannot complete on a regulation
9-ft table either. They are censored by the foot cushion, characterized
via the closed form, and the design was never re-aimed (a registered
exclusion, not a waiver).

\subsection{Group D --- curved paths (Coriolis), per experiment}
\label{sec:supp-results-d}

\paragraph{Parabolic slide path (\expid{D1})}
Experiment \expid{D1} verifies that the cue ball's sliding path after an
oblique rolling-ball impact is parabolic: half-ball hits at 1, 2, and
3\,m/s, both cut signs, sampled at 1\,ms. The registered acceptance
criterion --- quadratic-fit $R^2 > 1 - 10^{-6}$ in the frame aligned with
the measured post-impact velocity --- fails on all six runs, with
$1 - R^2$ between $2.2\times10^{-6}$ and $3.7\times10^{-6}$. The miss is a
frame artifact, and we keep it as this paper's worked example of a
frame-sensitive verification criterion: the real-ball impact model ($e_b = 0.95$,
finite ball--ball friction) leaves the friction acceleration
$2.9$--$4.1^\circ$ off perpendicular to that axis, injecting a cubic term
that a quadratic fit cannot absorb. In the slip-aligned frame --- the
parabola's own axis frame --- the fit is exact to machine precision:
$1 - R^2 = 0.0$ on every run, worst time-domain residual
$3.6\times10^{-16}$\,m, and the TP A.4 Eq.~20 slide duration reproduced to
within $2.8\times10^{-17}$\,s \citep{alciatore_tp_catalog}
(Fig.~\ref{fig:F10}a). The path extent scales as $v^{2.09}$--$v^{2.14}$
across the speed grid against the registered $\sim v^2$ claim, and the
final direction is nearly speed-invariant: the comparison error against
the TP A.4 closed-form angle anchor is $+0.17^\circ$ at 1\,m/s (the
post-cascade value) growing to $-0.94^\circ$ at 3\,m/s, a $1.11^\circ$
spread over the grid.

\paragraph{Mass\'e aiming (\expid{D2})}
Experiment \expid{D2} tests the mass\'e BAR aiming rule of
TP A.19 \citep{alciatore_tp_catalog} --- ``B for ball contact point, A
for cloth aim point, R for CB resting point,'' the source's own naming
\citep{alciatore_faq}: the final cue-ball direction should
parallel the line from the resting point (R) through the aim point (A)
on the cloth. Over elevations $45$--$75^\circ$ (a declared executed subset of
the registered $45$--$80^\circ$ range), eight azimuths, and two tip
speeds (64 cells; two ill-conditioned cells flagged and excluded from the
headline statistics), the final-heading error has mean $0.68^\circ$ and
maximum $1.60^\circ$, improving with elevation to a maximum of
$0.49^\circ$ at the grounded $75^\circ$ full-pinch mass\'e (the
$45$--$55^\circ$ cells run in the airborne jump-launch regime). Two
properties stand out. First, Coriolis' invariance theorem --- friction
acting along the fixed slip direction cannot rotate the final direction
\citep{coriolis1835} --- is realized exactly: the per-cell heading errors
at the two tip speeds are identical to the full recorded precision of the
output in all 32 (elevation, azimuth) pairs, so the per-speed error
statistics coincide bit-for-bit (Fig.~\ref{fig:F10}b). Second, 88\% of
cells ``come up a little short'' of the predicted direction, which is
precisely the behavior the source itself expects of real shots (TP A.19,
p.~3) --- a pass with the source's own caveat reproduced. The rule is also
usable where it is taught: the final direction is realized within 0.75\,m
at worst (median 0.18\,m), never beyond half a table length.

A scope note on the sub-pinch rows. The $45$--$55^\circ$ cells run in
the airborne jump-launch regime (the elevated strike launches the ball
off the cloth), and there everything \emph{after} the strike is inert to
the comparison by a conservation identity: under any impulsive strike
the natural-roll final heading --- the azimuth of the table-plane
combination $\mathbf{w}_{xy} = (5\mathbf{v}_{0,xy} +
2R\,\boldsymbol{\omega}_0 \times \hat{\mathbf{z}})/7$, the
$\times\hat{\mathbf{z}}$ term planar by construction --- is exactly
invariant to intervening hops, to the ball--slate restitution, and to
the sliding-friction law: every cloth-contact friction impulse enters
$\mathbf{v}_{xy}$ and $R\,\boldsymbol{\omega}\times\hat{\mathbf{z}}$
with equal and opposite weights in that combination, and a normal
(restitution) impulse at the cloth contact carries no planar component
and no torque about the center. Coriolis makes the robustness claim
himself --- the constructions giving the final directions ``subsistent
toujours, quelles que soient les lois que suivent les deux frottemens
de glissemens et de roulement'' \citep[quoted in][]{bouleau2006coriolis}
--- and the property is proven independently for this kernel's
strike--flight--contact pipeline (a launch-invariance theorem in the
kernel's unpublished design record). The $45$--$55^\circ$ rows therefore exercise
the strike's impulse--momentum mapping together with the flight and
contact solver's conservation structure --- a strike defect that
mis-mapped tip offset to launch velocity or spin would still miss the
R$\to$A line, so the rows are not vacuous --- but within the
impulse--momentum model class the BAR construction is reproduced by
derivation, so agreement there cannot discriminate strike \emph{realism}
beyond that idealization (finite tip dwell, tip--slate jamming) at those
elevations; we also know of no published mass\'e-aiming measurement in
the $30$--$60^\circ$ elevation band. The flight-leg half of the claim is
directly falsifiable in-sim: perturbing the ball--slate restitution
$e_t$ at $45^\circ$ elevation must leave each cell's BAR heading error
unchanged to output precision --- any larger shift would indicate a
broken flight/contact solver, not new strike physics.

\paragraph{Swerve maps (\expid{D3})}
Experiment \expid{D3} maps swerve over elevations
$\{0, 5, 10, 20, 30\}^\circ$, english fractions $\{25, 50, 100\%\}$, and
tip speeds 1--3\,m/s, with drift read at 0.5--2.5\,m along the aim line
(TP B.23 provenance; the registered claims are the FAQ's structural rules,
\citealp{alciatore_tp_catalog, alciatore_faq}). All three registered
structural rules hold with zero violations on the grid: no drift at level
cue (the maximum $|$curve$|$ over all level-cue rows is exactly zero ---
pure vertical-axis spin leaves the sliding leg straight), drift monotone
in elevation and in spin, and curvature at fixed distance decreasing with
speed, the delayed-onset rule. The delayed onset is made quantitative
through the squirt--swerve cancellation distance (``squerve''): the
distance at which swerve re-crosses the aim line that the $1.9^\circ$
full-english squirt displaced. Measured cancellation distances span
$0.08$--$2.65$\,m on the grid; at $5^\circ$ elevation and full english the
distance moves outward from 0.26\,m at 1\,m/s to 1.05\,m at 2\,m/s (the
marked $5^\circ$ crossing in Fig.~\ref{fig:F10}c) to
2.36\,m at 3\,m/s --- a derived, teachable quantity
linking the stick--ball and curved-path channels. Cells in which the ball
rests, or the 3.2\,m observation window ends, before re-crossing record no
distance. The sliding legs are closed-form parabolas to within
$2.6\times10^{-13}$\,m --- the same verification floor as \expid{D1} ---
and the experiment is cue-only, bit-identical across the slip-law kernel
cascade.

% ===========================================================================
% Supplement — full predicted-shift verification table (moved from the main
% text's Section 6 when the main text was shortened; the main text keeps
% a two-row extract, tab:predicted-shifts). Every number traces to the
% drift rows of the study's baseline manifest and its ledger under the
% canonical kernel 0.1.0+fb3425ad.
% ===========================================================================
\section{Predicted-shift verification across the slip-law and
  branch-selection fix cascade}
\label{sec:supp-predicted-shifts}

Before the post-fix battery re-run, each affected experiment received a
written prediction of the direction and mechanism of its shift, and no
registered band or catalog parameter was re-tuned
(Section~\ref{sec:oracle} of the main text presents the fix-cascade case
study and a two-row extract). Table~\ref{tab:supp-predicted-shifts}
transcribes all seven prediction classes; all seven confirmed.

% Table typography: the longtable takes the
% sanctioned 7.0 in artwork width via a symmetric \LTleft/\LTright
% overhang (matching the native-size figure convention; 'plus 1fill'
% keeps the chunks centered without per-row hbox warnings), with
% absolute ragged-right column widths and continuation furniture. The
% col-1 header is abbreviated 'Expt.' so it no longer overflows into
% col 2 (row labels C1/A3/... are defined in the text).
{\small
\setlength{\LTleft}{-57pt plus 1fill}%
\setlength{\LTright}{-57pt plus 1fill}%
\begin{longtable}{@{}>{\raggedright\arraybackslash}p{0.55in}>{\raggedright\arraybackslash}p{2.9in}>{\raggedright\arraybackslash}p{3.19in}@{}}
  \caption{Predicted-shift verification across the slip-law and
  branch-selection fix cascade.
  Each row is a prediction class registered before the post-fix battery
  re-run; ``bit-identical'' means byte equality of the affected rows
  after masking kernel-version stamps. No acceptance band or catalog
  parameter was re-tuned.}
  \label{tab:supp-predicted-shifts}\\
    \toprule
    Expt. & Registered prediction & Measured outcome \\
    \midrule
    \endfirsthead
    \multicolumn{3}{@{}l}{Table~\ref{tab:supp-predicted-shifts} \emph{(continued)}}\\
    \toprule
    Expt. & Registered prediction & Measured outcome \\
    \midrule
    \endhead
    \midrule
    \multicolumn{3}{r@{}}{\emph{Continued on next page}}\\
    \endfoot
    \bottomrule
    \endlastfoot
    \expid{C1} &
    branch-selection fix: the four \code{stronge_realistic} discontinuous
    bank cells collapse;
    other presets untouched &
    cells $4\to0$; the three-crossing cell (aim 3.0, 1\,m/s) monotone;
    \code{han2005}, \code{mathavan}, and \code{grip} rows bit-identical \\
    \addlinespace
    \expid{A3} &
    slip tan-argument ratio $\to 1$ (throw exactly $\arctan\mu$); stick
    regime unchanged &
    ratio $0.9745\to1.00014$; stick ratio 1.02564 unchanged; worst
    slip-channel $|E|$ vs the transcribed master curves
    $0.0886\to0.0012^\circ$ \\
    \addlinespace
    \expid{A1} &
    half-ball separations shift in the $+2.6$\% slip-throw direction;
    straight-in rows unchanged &
    separations $-0.10$ to $-0.13^\circ$; straight-in rows
    bit-identical \\
    \addlinespace
    \expid{A5}, \expid{A6} &
    gearing zero-crossings bit-identical (gearing zeroes slip);
    slip-branch ratios rescale by $(1+e_b)/(2e_b)$ &
    zero-crossings bit-identical; transfer ratio
    $0.95007\to0.97508 = (1+e_b)/2$; half-ball sanity throw
    $4.5222\to4.6407^\circ$ \\
    \addlinespace
    \expid{B1}--\expid{B4} &
    strike-only experiments unchanged &
    \expid{B1}, \expid{B2}, \expid{B4} bit-identical; \expid{B3}
    identical except a stale script stamp flag (bookkeeping) \\
    \addlinespace
    \expid{D1}--\expid{D3}, \expid{F1}--\expid{F3} &
    \expid{D2}, \expid{D3}, \expid{F1} unchanged; \expid{F3} shifts only
    in its over-cut family (airborne contacts are slip-branch impacts);
    \expid{D1} shifts (post-impact slide is slip-branch) &
    confirmed as stated: \expid{D1} asymptote $E_\theta$ at 1\,m/s
    $0.1307\to0.1729^\circ$; \expid{F3} min-elevation curves and jump
    sweep bit-identical \\
    \addlinespace
    \expid{E1}, \expid{E3} &
    compliant-cluster rows unchanged (force-based Coulomb consumes no
    impulse cap); the model contrast survives &
    \expid{E1} compliant slope bit-identical (contrast 0.139 vs 0.0295);
    \expid{E3} compliant squat rows bit-identical \\
\end{longtable}
}

% ===========================================================================
% Supplement S6 — presentation methods for the validated-scenario video
% artifacts: the C² log-domain time remap and shutter synthesis. Sources
% of truth: the shipped implementation (web/app/framestep.js —
% smootherstep / logRemapFrames / shutterSubframes, the one frame-timeline
% module shared by the in-app exporter and the node render harness) and
% the cited primary sources (entry-level metadata verified against them on
% 2026-07-26; see references.bib).
% SCOPE GUARD (taxonomy discipline): the simulator's
% output of record is the kernel event stream (Section 3 of the main
% text); this section documents how gated records are PRESENTED as video,
% and nothing in it enters the paper's V&V taxonomy or evidence chain.
% ===========================================================================
\section{Presentation methods for the validated-scenario video artifacts}
\label{sec:supp-timeremap}

The artifact bundle's validated-scenario video pack
(Section~\ref{sec:appendix-repro}) renders sixteen catalog entries ---
each a solved strike, or a short ladder of strikes, replayed through the
canonical binary shipped at the time of its render --- as annotated
slow-motion videos. The simulator's output of record is the
kernel event stream, not any sampled rendering of it (main text,
Section~3), and no claim of the paper rests on these videos: each entry
first re-simulates its claimed scenario on that shipped canonical binary
and must reproduce its pinned record digest \emph{before} any frame is
captured, so the videos are presentations of already-verified records.
The per-entry kernel stamp is recorded in the entry's gate report; the
pack's stamps postdate the paper's battery pin along the supersession
lineage of the main text, and each entry's digests are pinned against
its own stamp.
This section documents the deterministic time-remap and shutter
synthesis used by those presentations, for two reasons: the frame
schedule is deterministic --- a pure function of the entry's playback
plan, which is authored in the render tool's scenario catalog, with the
derived schedule's summary (frame counts, worst rate-step ratio,
shutter parameters, replay spans) pinned in the per-entry gate
report --- and the design choices are
non-obvious enough that a reproducer would otherwise have to
reverse-engineer them from the delivered bytes.

\paragraph{Frame schedule and the log-domain speed ramp.}
Videos are delivered at a constant 60 frames per second; one output
frame advances the kernel clock by $\Delta = 1/60$~s times a
dimensionless playback rate $r$ (kernel seconds consumed per output
second; $r = 1$ is real time). A playback plan is a short list of rate
segments over kernel time, each pair of consecutive segments separated
either by a camera cut (a hard rate step on the cut frame) or by a
\emph{ramp}. A hard rate step without a cut reads as a speed jerk, so
cut-less boundaries between rates $r_1$ and $r_2$ are ramped in the
logarithmic rate domain:
\begin{equation}
r(u) \;=\; r_1\,\bigl(r_2/r_1\bigr)^{S(u)},
\qquad
S(x) \;=\; x^3\,(6x^2 - 15x + 10),
\qquad u \in [0,1],
\label{eq:supp-lograte}
\end{equation}
where $S$ is the quintic ``smootherstep'' polynomial, with zero first
and second derivatives at both endpoints; the remap is therefore $C^2$
in log-rate, so the on-screen speed carries no visible acceleration
kink. $S$ is, term for term, the minimum-jerk rest-to-rest trajectory of
human reaching movements \citep[eq.~(2), p.~1689]{flash1985coordination}:
shaping the ramp with it makes the log-rate trajectory the minimum-jerk
path between $\log_2 r_1$ and $\log_2 r_2$. The ramp duration follows a
per-octave law,
\begin{equation}
D \;=\; \operatorname{clamp}\bigl(0.30 \cdot \lvert \log_2 (r_2/r_1)
\rvert,\; 0.5,\; 1.5\bigr)~\text{output seconds},
\label{eq:supp-octave}
\end{equation}
quantized to $n = \max(1, \operatorname{round}(D/\Delta))$ whole output
frames --- a fixed time cost per doubling of speed, bounded below and
above so that small steps are not abrupt and large ones do not drag.
The per-frame kernel advance integrates the ramp at the frame midpoint,
$\mathrm{d}t_j = \Delta \cdot r((j + \tfrac12)/n)$ for
$j = 0,\dots,n-1$, and the ramp is scheduled as a \emph{lead-in}: it
starts on the last steady frame from which its integrated kernel
consumption still completes at or before the segment boundary, so that
whenever the outgoing segment can host the ramp --- the common case ---
the target rate is established at or before the boundary and impacts
land fully inside their slow-motion rate. A segment too short to host
its ramp instead starts the ramp at segment entry and overshoots the
boundary; the published pack contains one such replay pass, and the
exception is why the rate-step law below and the per-entry gate windows
are measured on the \emph{emitted} clock, never on the nominal plan.
Plan validation rejects
non-positive or non-finite rates, which the log-domain form requires
(the analogous restriction in the Web Audio specification is discussed
below).

\paragraph{Perceptual basis for ramping in the log domain.}
The ramp is built in the logarithmic speed domain, where equal steps are
equal speed \emph{ratios}; the quintic easing modulates how fast that
axis is traversed (its rate reaching zero at both endpoints), so the
constant-ratio character lives in the easing variable $S(u)$ rather
than literally in each output frame. The ratio axis is the perceptually
relevant one for visual speed. The classical measurement is Weber-law
behavior: the
just-noticeable difference in speed grows in proportion to the reference
speed \citep{brown1960weber}. \citet{debruyn1988velocity} measured the
Weber fraction for velocity discrimination at approximately 7\% at the
optimal speeds (84\%-correct staircase criterion), roughly constant over
a 4--64~deg/s plateau and rising (U-shaped) at slower and faster
speeds. In modeling terms, \citet{stocker2006noise} adopt a logarithmic
speed representation for their likelihood and internal-noise functions
(a log-Normal likelihood, $\tilde{v} = \ln((v + v_0)/v_0)$, in their
Methods), motivated by exactly this Weber--Fechner character of speed
discrimination --- a modeling choice for the noise function, not a
direct neural-coding proof; the neural grounding of a logarithmic,
scale-invariant speed representation is the middle-temporal-area
population analysis of \citet{nover2005logarithmic}. A perceptually
equal speed change is therefore an equal step in $\log r$, and
Eq.~\eqref{eq:supp-lograte} ramps along exactly that axis, its
minimum-jerk easing distributing the steps in time. The same reasoning is
codified for audio in the W3C Web Audio API, whose
\code{exponentialRampToValueAtTime} automation has the identical
functional form $v(t) = V_0 (V_1/V_0)^{(t - T_0)/(T_1 - T_0)}$ and whose
specification recommends exponential ramps for ``filter frequencies and
playback rate . . . because of the way humans perceive sound''
\citep{webaudio_spec}. The specification's domain restriction is
instructive: an exponential ramp can neither reach nor pass through
zero --- the W3C text throws a \code{RangeError} only when the target
value equals zero, while a zero or opposite-signed \emph{start} value
silently freezes the ramp (a ramp between two negative values is
permitted). Our plan validator is intentionally stricter than that
domain: playback rates must be finite and strictly positive, so the
degenerate cases the specification handles by freezing or throwing
cannot be expressed in a plan at all.

\paragraph{Shutter synthesis over world-time windows.}
Motion blur is synthesized, not post-processed: each delivered frame is
the accumulation of 8 sub-frame renders spanning a fixed
$310^\circ$ shutter angle, i.e.\ the shutter opens at the frame's kernel
time $t_j$ and stays open for $(310/360) \cdot \mathrm{d}t_j$ of
\emph{kernel} (world) time, where $\mathrm{d}t_j$ is the integrated
per-frame advance above --- never the nominal segment rate. Blur length
therefore follows the on-screen speed through ramps exactly as a
physical camera filming the slow-motion playback would behave. In the
frame domain the sub-frame animation steps sum to exactly $\Delta$ per
output frame (the first sub-frame carries the closed-shutter
remainder), so frame counts and presentation-layer animation lengths
are independent of the shutter setting.

\paragraph{Monotone-map anchoring.}
The schedule is emitted in monotone runs: within a run the map from
output frame index to kernel time is strictly increasing, and a
broadcast-style replay (a rewind to re-cover an impact) starts a new
run, its boundary marked by a camera cut and a 0.4~s cross-dissolve
identifying re-covered time. The \emph{placement} of timed presentation
content ---
predicted-overlay geometry, measured-value callouts, story-beat text
reveals --- is declared as \emph{kernel-time} windows and resolved to
output frames through this monotone map, never positioned by hand in
output seconds; within a window, the reveal choreography itself (alpha
fades, type entrances) plays authored fixed output-time durations ---
0.4~s fades matching the dissolve above --- and a few entries carry
small authored output-time leads. An annotation window that would match
more than one run of its
camera is rejected rather than guessed at. One mechanical law bounds
the emitted clock: the worst consecutive cut-less per-frame
kernel-advance ratio, measured on the emitted clock per monotone run,
must not exceed $1.35\times$; the measured value is recorded in the
per-entry gate report, and an entry that exceeds the bound fails before
any frame is captured.

\paragraph{Relation to prior art.}
The two ingredients of Eq.~\eqref{eq:supp-lograte} appear separately in
published work. Speed-template curves in video-editing tools enumerate
logarithmic and exponential playback-speed transitions among other
shapes, motivated by editorial smoothness rather than psychophysics
\citep{gopro2017variable,adobe2012changing}; and Weber--Fechner-guided
playback-rate control appears in adaptive fast-forward summarization,
where perceived fast-forward speed is modeled on a logarithmic axis
\citep{jin2017elasticplay}. We have not found a published description
that combines a perceptually motivated log-domain speed curve with
$C^2$ minimum-jerk easing for cinematic slow-motion time remapping; we
record the synthesis here as presentation methodology --- it is not a
claim of this paper, and none of the paper's results depend on it.


% ---------------------------------------------------------------------------
\bibliographystyle{elsarticle-harv}
\bibliography{references}

\end{document}
